%%%%%%%%%%%%%%%%%%%%%%%%%%%%%%%%%%%%%%%%%%%%%%%%%%%%%%%%%%%%%%%%%%%%%%%%%%%%%%%
%  SÍLABO DE MACROECONOMÍA — Plan de estudio integrado
%
%  Textos base:
%    DG — José De Gregorio, Macroeconomía: Teoría y Políticas (Pearson, 2007)
%    H  — Ben J. Heijdra, Foundations of Modern Macroeconomics (OUP, 3.ª ed., 2017)
%    R  — David Romer, Advanced Macroeconomics (McGraw-Hill, 5.ª ed., 2019)
%    CSV — Campante, Sturzenegger y Velasco, Advanced Macroeconomics:
%          An Easy Guide (LSE Press, 2021)
%    A  — Daron Acemoglu, Introduction to Modern Economic Growth (Princeton, 2009)
%    W  — Carl E. Walsh, Monetary Theory and Policy (MIT Press, 3.ª ed., 2010)
%    G  — Jordi Galí, Monetary Policy, Inflation, and the Business Cycle
%          (Princeton, 2008)
%  Consulta:
%    GM — Galindo y Montecinos, Macrodinámica, Tomo I: Modelos de Ciclos
%          Económicos Reales (2017)
%  Suplemento:
%    K  — Azzimonti, Krusell, McKay y Mukoyama, Macroeconomics (OUP, 2026)
%  Programas PUCP integrados: 1ECO74, 2ECO03, 1ECO90, 1ECO14, ECO324, ECO338,
%  MAT616
%
%  Compilar con:  pdflatex Macroeconomia.tex   (dos veces, por el índice)
%%%%%%%%%%%%%%%%%%%%%%%%%%%%%%%%%%%%%%%%%%%%%%%%%%%%%%%%%%%%%%%%%%%%%%%%%%%%%%%

\documentclass[11pt,a4paper]{article}

\usepackage[T1]{fontenc}
\usepackage[utf8]{inputenc}
\usepackage[spanish,es-noquoting,es-tabla]{babel}
\usepackage{lmodern}
\usepackage[scaled=0.92]{helvet}        % títulos en sans serif, al estilo McGraw-Hill
\usepackage[margin=2.4cm,headheight=15pt]{geometry}
\usepackage[table]{xcolor}
\usepackage{enumitem}
\usepackage{titlesec}
\usepackage{titletoc}
\usepackage{tabularx}
\usepackage{booktabs}
\usepackage{longtable}
\usepackage{amsmath}
\usepackage{amssymb}
\usepackage{fancyhdr}
\usepackage{tikz}
\usepackage[skins]{tcolorbox}
\usepackage{needspace}

%------------------------------------------------------------------ colores ---
% Gama de fucsias y magentas
\definecolor{tonoOscuro}{RGB}{82,4,54}              % oscuro
\definecolor{tonoProfundo}{RGB}{121,6,79}           % profundo
\definecolor{tonoPrincipal}{RGB}{160,8,104}         % principal
\definecolor{tonoA}{RGB}{179,25,92}                 % variante cálida
\definecolor{tonoB}{RGB}{179,15,146}                % variante fría
\definecolor{tonoC}{RGB}{193,21,187}                % variante viva
\definecolor{tonoClaro}{RGB}{221,95,175}            % claro
\definecolor{tonoAcero}{RGB}{131,63,106}            % acero
\definecolor{tonoPizarra}{RGB}{134,86,116}          % pizarra
\definecolor{tonoGrafito}{RGB}{73,39,61}            % grafito (texto secundario)
\definecolor{tonoHielo}{RGB}{247,207,232}           % hielo
\definecolor{tonoPalido}{RGB}{251,233,245}          % pálido
\definecolor{tonoBruma}{RGB}{254,246,251}           % bruma

\colorlet{azulosc}{tonoOscuro}
\colorlet{grisclaro}{tonoPalido}
% Un tono distinto de la gama para cada libro
\definecolor{colorDG}{RGB}{138,15,138}
\definecolor{colorH}{RGB}{184,20,169}
\definecolor{colorK}{RGB}{161,18,134}
\definecolor{colorR}{RGB}{138,15,104}
\definecolor{colorCSV}{RGB}{184,20,124}
\definecolor{colorGM}{RGB}{161,18,95}
\definecolor{colorA}{RGB}{138,15,70}
\definecolor{colorW}{RGB}{184,20,79}
\definecolor{colorG}{RGB}{161,18,56}
\colorlet{colorPUCP}{tonoPizarra}

\usepackage[colorlinks=true,linkcolor=tonoOscuro,urlcolor=tonoPrincipal]{hyperref}

%------------------------------------------------------------------ estilos ---
\titleformat{\section}[block]
  {\sffamily\LARGE\bfseries\color{tonoOscuro}\raggedright}{\thesection.}{0.6em}
  {\textcolor{tonoPrincipal}{\rule[-0.1ex]{0.55em}{0.85em}}\hspace{0.45em}}
  [\vspace{-0.4ex}{\color{tonoPrincipal}\titlerule[1.6pt]}\vspace{1.2pt}{\color{tonoClaro}\titlerule[0.6pt]}]
\titlespacing*{\section}{0pt}{2.6ex plus 1ex minus .2ex}{1.6ex plus .2ex}

\titleformat{\subsection}
  {\sffamily\large\bfseries\color{tonoPrincipal}}{\thesubsection}{0.6em}
  {\textcolor{tonoClaro}{\rule[0.1ex]{0.45em}{0.45em}}\hspace{0.45em}}
\titlespacing*{\subsection}{0pt}{2.2ex plus .8ex minus .2ex}{1ex plus .2ex}

% Índice
\titlecontents{section}[0pt]
  {\addvspace{3pt}\sffamily\bfseries\color{tonoOscuro}}
  {}{\textcolor{tonoB}{\rule[0.15ex]{0.4em}{0.4em}}\hspace{0.6em}}
  {{\color{tonoClaro}\titlerule*[0.7em]{.}}\textcolor{tonoPrincipal}{\contentspage}}

% Cabecera y pie
\pagestyle{fancy}
\fancyhf{}
\lhead{\small\sffamily\bfseries\color{tonoOscuro}Sílabo de Macroeconomía}
\rhead{\footnotesize\sffamily\color{tonoPrincipal}De Gregorio $\cdot$ Heijdra $\cdot$ Romer $\cdot$ Campante $\cdot$ Acemoglu $\cdot$ Walsh $\cdot$ Galí $\cdot$ Krusell}
\cfoot{{\setlength{\fboxsep}{3pt}\colorbox{tonoPrincipal}{\makebox[1.8em]{\small\sffamily\bfseries\color{white}\thepage}}}}
\renewcommand{\headrule}{%
  {\color{tonoPrincipal}\hrule height 1.1pt width\headwidth}\vskip 1.2pt
  {\color{tonoClaro}\hrule height 0.4pt width\headwidth}\vskip -2.7pt}

% Listas
\setlist[itemize,1]{label=\textcolor{tonoPrincipal}{\scriptsize$\blacksquare$}}
\setlist[enumerate,1]{label=\textcolor{tonoPrincipal}{\sffamily\bfseries\arabic*.}}

% Tablas: cabecera oscura y filas alternas en tono pálido
\arrayrulecolor{tonoPrincipal}
\setlength{\aboverulesep}{0pt}
\setlength{\belowrulesep}{0pt}
\renewcommand{\arraystretch}{1.3}
\newcommand{\cab}[1]{\textcolor{white}{\sffamily\bfseries #1}}
\newcommand{\filasalternas}{\rowcolors{2}{white}{tonoPalido}}

% Encabezados de párrafo en color
\newcommand{\destacado}[1]{{\sffamily\bfseries\color{tonoPrincipal}#1}}

% Franja decorativa con la gama
\newcommand{\franjas}[1][3pt]{%
  \noindent\begin{tikzpicture}
    \foreach \c [count=\i from 0] in {tonoOscuro,tonoProfundo,tonoPrincipal,tonoA,tonoB,tonoC,tonoClaro,tonoHielo}
      \fill[\c] (\i*\linewidth/8,0) rectangle ++(\linewidth/8,#1);
  \end{tikzpicture}\par}

%---------------------------------------------------------------- comandos ----
% Etiqueta de libro: recuadro de color con la sigla
\newcommand{\etiqueta}[2]{%
  {\setlength{\fboxsep}{1.4pt}\colorbox{#1}{\sffamily\bfseries\color{white}\strut#2}}}
% Referencias abreviadas a cada libro
\newcommand{\DG}[1]{\etiqueta{colorDG}{DG}\nobreak\hspace{0.25em}\textcolor{colorDG}{#1}}
\newcommand{\HJ}[1]{\etiqueta{colorH}{H}\nobreak\hspace{0.25em}\textcolor{colorH}{#1}}
\newcommand{\KR}[1]{\etiqueta{colorK}{K}\nobreak\hspace{0.25em}\textcolor{colorK}{#1}}
\newcommand{\RO}[1]{\etiqueta{colorR}{R}\nobreak\hspace{0.25em}\textcolor{colorR}{#1}}
\newcommand{\CSV}[1]{\etiqueta{colorCSV}{CSV}\nobreak\hspace{0.25em}\textcolor{colorCSV}{#1}}
\newcommand{\AC}[1]{\etiqueta{colorA}{A}\nobreak\hspace{0.25em}\textcolor{colorA}{#1}}
\newcommand{\WA}[1]{\etiqueta{colorW}{W}\nobreak\hspace{0.25em}\textcolor{colorW}{#1}}
\newcommand{\GA}[1]{\etiqueta{colorG}{G}\nobreak\hspace{0.25em}\textcolor{colorG}{#1}}
\newcommand{\GM}[1]{\etiqueta{colorGM}{GM}\nobreak\hspace{0.25em}\textcolor{colorGM}{#1}}

% Un tema: número en recuadro de color + título
\newcommand{\tema}[1]{\temaaux#1\relax}
\def\temaaux#1 #2\relax{%
  \par\addvspace{1.1em}\noindent
  {\sffamily\bfseries\large\raggedright
   \hangindent=2.9em\hangafter=1
   \makebox[2.9em][l]{{\setlength{\fboxsep}{2.5pt}\colorbox{tonoPrincipal}{\color{white}#1}}}%
   {\color{tonoOscuro}#2}\par}%
  \nobreak\vspace{0.25em}\noindent{\color{tonoClaro}\rule{\linewidth}{0.6pt}}\par\nobreak\vspace{0.35em}}

% Línea de lecturas con marcador y etiqueta de color
\newcommand{\linealectura}[3]{%
  \par\nobreak\noindent\hspace{1.2em}%
  \begin{minipage}{0.92\textwidth}\small
    \textcolor{#1}{\rule[0.1ex]{0.45em}{0.45em}}\hspace{0.45em}%
    {\sffamily\bfseries\color{#1}#2}\enspace#3%
  \end{minipage}\par\vspace{0.3em}}
\newcommand{\lecturas}[1]{\linealectura{tonoPrincipal}{Lecturas:}{#1}}
\newcommand{\apoyo}[1]{\linealectura{tonoA}{Profundización:}{#1}}
\newcommand{\juegos}[1]{\linealectura{tonoB}{Teoría de juegos:}{#1}}
\newcommand{\supl}[1]{\linealectura{tonoC}{Suplemento:}{#1}}
\newtcolorbox{cajanota}{enhanced, sharp corners, boxrule=0pt, frame hidden,
  colback=tonoPalido, borderline west={2.5pt}{0pt}{tonoClaro},
  left=8pt, right=6pt, top=4pt, bottom=4pt, before skip=0pt, after skip=0pt,
  fontupper=\small\color{tonoGrafito}}
\newcommand{\nota}[1]{%
  \par\nobreak\noindent\hspace{1.2em}%
  \begin{minipage}{0.92\textwidth}\begin{cajanota}#1\end{cajanota}\end{minipage}\par\vspace{0.3em}}
\newcommand{\consultalib}[1]{\linealectura{tonoPizarra}{Consulta:}{{\color{tonoGrafito}#1}}}
% Artículos y textos tomados de los programas PUCP
\newcommand{\articulos}[1]{\linealectura{colorPUCP}{Artículos (PUCP):}{#1}}
% Artículos que no proceden de los programas PUCP (entran en las horas)
\newcommand{\articulosB}[1]{\linealectura{tonoB}{Artículos:}{#1}}
\newcommand{\consulta}[1]{\linealectura{colorPUCP}{Consulta (PUCP):}{{\color{tonoGrafito}#1}}}

% Encabezado de bloque: banda oscura, franja de la gama y recuadro de presentación
\newif\ifrecienparte
\newcommand{\bloque}[2]{%
  \ifrecienparte\global\recienpartefalse\else\par\needspace{0.4\textheight}\fi
  \phantomsection\addcontentsline{toc}{section}{#1}%
  \begin{tcolorbox}[enhanced, sharp corners, boxrule=0pt, frame hidden,
      colback=tonoOscuro, borderline west={7pt}{0pt}{tonoClaro},
      left=14pt, right=10pt, top=10pt, bottom=10pt, before skip=0pt, after skip=0pt]
    {\raggedright\sffamily\Large\bfseries\color{white}#1\par}
  \end{tcolorbox}%
  \nointerlineskip\franjas[4pt]
  \vspace{0.8em}
  \begin{tcolorbox}[enhanced, sharp corners, boxrule=0pt, frame hidden,
      colback=tonoPalido, borderline west={3pt}{0pt}{tonoPrincipal},
      borderline south={0.6pt}{0pt}{tonoClaro},
      left=10pt, right=8pt, top=6pt, bottom=6pt, before skip=0pt, after skip=0pt,
      fontupper=\small\color{tonoGrafito}]
    #2
  \end{tcolorbox}%
  \par\medskip}

%================================================================== documento ==
\begin{document}

\begin{titlepage}
\begin{tikzpicture}[remember picture,overlay]
  % banda superior
  \fill[tonoOscuro] (current page.north west) rectangle ([yshift=-5.4cm]current page.north east);
  \fill[tonoClaro] (current page.north west) rectangle ([xshift=1.1cm,yshift=-5.4cm]current page.north west);
  % franja de la gama bajo la banda
  \foreach \c [count=\i from 0] in {tonoOscuro,tonoProfundo,tonoPrincipal,tonoA,tonoB,tonoC,tonoClaro,tonoHielo}
    \fill[\c] ([xshift=\i*\paperwidth/8,yshift=-5.4cm]current page.north west)
      rectangle ++(\paperwidth/8,-0.35cm);
  % franja inferior
  \foreach \c [count=\i from 0] in {tonoHielo,tonoClaro,tonoC,tonoB,tonoA,tonoPrincipal,tonoProfundo,tonoOscuro}
    \fill[\c] ([xshift=\i*\paperwidth/8]current page.south west)
      rectangle ++(\paperwidth/8,0.6cm);
\end{tikzpicture}%
\centering
\vspace*{0.5cm}
{\sffamily\Huge\bfseries\color{white} Sílabo de Macroeconomía\par}
\vspace{0.5cm}
{\sffamily\Large\color{tonoHielo} Plan de estudio integrado\par}
\vspace{1.9cm}
{\sffamily\large\bfseries\color{tonoPrincipal}\textcolor{tonoClaro}{\rule[0.5ex]{1.5cm}{0.6pt}}\quad Textos base\quad\textcolor{tonoClaro}{\rule[0.5ex]{1.5cm}{0.6pt}}\par}
\vspace{0.35cm}
{\large\etiqueta{colorDG}{DG} \quad {\color{tonoOscuro}José De Gregorio}\par}
{\normalsize \emph{\color{colorDG}Macroeconomía: Teoría y Políticas}\par}
{\small\color{tonoPizarra} Pearson Educación, México, 2007 --- 25 capítulos\par}
\vspace{0.3cm}
{\large\etiqueta{colorH}{H} \quad {\color{tonoOscuro}Ben J. Heijdra}\par}
{\normalsize \emph{\color{colorH}Foundations of Modern Macroeconomics}\par}
{\small\color{tonoPizarra} Oxford University Press, 3.ª edición, 2017 --- 19 capítulos\par}
\vspace{0.3cm}
{\large\etiqueta{colorR}{R} \quad {\color{tonoOscuro}David Romer}\par}
{\normalsize \emph{\color{colorR}Advanced Macroeconomics}\par}
{\small\color{tonoPizarra} McGraw-Hill Education, 5.ª edición, 2019 --- 13 capítulos\par}
\vspace{0.3cm}
{\large\etiqueta{colorCSV}{CSV} \quad {\color{tonoOscuro}Campante, Sturzenegger y Velasco}\par}
{\normalsize \emph{\color{colorCSV}Advanced Macroeconomics: An Easy Guide}\par}
{\small\color{tonoPizarra} LSE Press, 2021 --- 22 capítulos\par}
\vspace{0.6cm}
{\sffamily\large\bfseries\color{tonoPrincipal}\textcolor{tonoClaro}{\rule[0.5ex]{1.5cm}{0.6pt}}\quad Textos de referencia de posgrado\quad\textcolor{tonoClaro}{\rule[0.5ex]{1.5cm}{0.6pt}}\par}
\vspace{0.35cm}
{\large\etiqueta{colorA}{A} \quad {\color{tonoOscuro}Daron Acemoglu}\par}
{\normalsize \emph{\color{colorA}Introduction to Modern Economic Growth}\par}
{\small\color{tonoPizarra} Princeton University Press, 2009 --- 24 capítulos y 3 apéndices\par}
\vspace{0.3cm}
{\large\etiqueta{colorW}{W} \quad {\color{tonoOscuro}Carl E. Walsh}\par}
{\normalsize \emph{\color{colorW}Monetary Theory and Policy}\par}
{\small\color{tonoPizarra} MIT Press, 3.ª edición, 2010 --- 11 capítulos\par}
\vspace{0.3cm}
{\large\etiqueta{colorG}{G} \quad {\color{tonoOscuro}Jordi Galí}\par}
{\normalsize \emph{\color{colorG}Monetary Policy, Inflation, and the Business Cycle}\par}
{\small\color{tonoPizarra} Princeton University Press, 2008 --- 8 capítulos\par}
\vspace{0.6cm}
{\large\etiqueta{colorK}{K} \quad {\color{tonoOscuro}Azzimonti, Krusell, McKay y Mukoyama}\par}
{\normalsize \emph{\color{colorK}Macroeconomics}\par}
{\small\color{tonoPizarra} Oxford University Press, 2026 --- \textbf{texto suplementario}\par}
\vspace{0.35cm}
{\small\etiqueta{colorGM}{GM} \quad Galindo y Montecinos, \emph{Macrodinámica, Tomo I} (2017) --- \textbf{consulta}\par}
\vspace{0.4cm}
{\normalsize\color{colorPUCP} Con temas y lecturas de siete programas de la PUCP\par}
\vfill
{\setlength{\fboxsep}{6pt}\colorbox{tonoPrincipal}{\normalsize\sffamily\bfseries\color{white} Quince bloques temáticos --- unos 22 meses}\par}
\vspace{0.4cm}
\end{titlepage}

\tableofcontents
\newpage

%------------------------------------------------------------------------------
\section*{Cómo usar este sílabo}
\addcontentsline{toc}{section}{Cómo usar este sílabo}

El plan está organizado en \textbf{quince bloques temáticos}, no por capítulos de
un libro u otro. La numeración de los textos no se corresponde, de modo que
intentar leerlos en paralelo página a página resulta contraproducente. La
recomendación es cerrar cada bloque antes de abrir el siguiente, en el orden que
fija el cronograma: los Bloques 1 a 11 van en orden; luego se estudian los
métodos recursivos del Bloque 14 (temas 14.1 a 14.7), después el Bloque 12, que
intercala el tema 14.8, y al final los Bloques 13 y 15. La numeración de los
bloques se mantiene porque otros sílabos de la carpeta remiten a ella.

\medskip
\noindent\destacado{Criterio general de reparto} (en el orden en que se lee cada tema):
\begin{itemize}[leftmargin=1.4em,itemsep=2pt]
  \item \textcolor{colorDG}{\textbf{De Gregorio}} aporta la intuición económica,
        la evidencia empírica y el tratamiento de la economía abierta pequeña.
        Se lee \emph{primero}, para entender cuál es el problema.
  \item \textcolor{colorCSV}{\textbf{Campante, Sturzenegger y Velasco}} presentan
        cada modelo completo en pocas páginas, con mirada de economía emergente:
        son el puente hacia la formalización.
  \item \textcolor{colorH}{\textbf{Heijdra}} aporta la formalización, la dinámica
        y el puente hacia la literatura de investigación. Se lee \emph{después},
        para resolver el problema con rigor.
  \item \textcolor{colorR}{\textbf{Romer}}, texto de referencia de los programas
        de posgrado, es el segundo tratamiento formal: se lee con horas en los
        temas que Heijdra no cubre (evidencia empírica, contratos, economía
        política de los déficits, fricciones financieras) y como consulta en los
        que Heijdra ya resuelve.
  \item \textcolor{colorA}{\textbf{Acemoglu}}, \textcolor{colorW}{\textbf{Walsh}} y
        \textcolor{colorG}{\textbf{Galí}}, referencias de posgrado en crecimiento y
        en teoría y política monetaria, cierran los temas de su área y son la
        lectura principal de los nuevos: Acemoglu en los Bloques 4 a 6, Walsh en
        el 8 y el 9, y Galí en el 12 y el 13.
\end{itemize}

\medskip
\noindent\destacado{Un texto por nivel.} En los bloques de crecimiento (4 a 6) y en
el Bloque 1 cada tema se lee en un texto de intuición (De Gregorio o Campante et
al.), uno formal (Heijdra o, donde Heijdra no llega, Romer) y uno de
profundidad (Acemoglu). Los demás tratamientos del mismo modelo, incluido el
suplemento, quedan como consulta y no suman horas.

\medskip
\noindent\destacado{Material excluido.} Este plan arranca directamente en el
Bloque 1. Queda fuera el material introductorio de De Gregorio y Heijdra —panorama de
escuelas de pensamiento, cuentas nacionales y convenciones de medición
(\DG{1}, \DG{2}, \HJ{1})— que se da por conocido. Si en algún momento surge una
duda sobre definiciones contables, deflactores o balanza de pagos, \DG{2} es la
referencia a la que volver. Por la misma razón se omiten la introducción de
\textcolor{colorR}{\textbf{R}} y el capítulo 1 de \textcolor{colorCSV}{\textbf{CSV}}.

\medskip
\noindent\destacado{Sobre las matemáticas.} El plan supone el sílabo de economía
matemática, en particular su Parte II: ecuaciones diferenciales y en
diferencias, sistemas dinámicos y diagramas de fase (Bloques 9 y 10), control
óptimo (Bloque 12) y programación dinámica (Bloque 13). Conviene haber cerrado
los Bloques 9 y 10 de ese plan antes del tema 1.6, el Bloque 12 antes del
Bloque 4 y el Bloque 13 antes del Bloque 14 de este. Por eso este sílabo no
tiene pretramo matemático, y los apéndices matemáticos de Heijdra, Acemoglu y
Campante et al. quedan solo como consulta. La programación dinámica a nivel de
posgrado, la solución de modelos con expectativas racionales y el cómputo en
Python ocupan el Bloque 14, que se estudia antes del 12.

\medskip
\noindent\destacado{Sobre los ejercicios.} El \emph{Exercise and Solutions Manual}
de Heijdra es prácticamente obligatorio. Estos modelos no se aprenden leyendo,
se aprenden derivándolos con lápiz y papel. Los problemas de fin de capítulo de
Romer, Acemoglu y Walsh cumplen la misma función; el Anexo G reúne
material opcional de práctica, de consulta y de macroeconomía e inteligencia artificial.

\medskip
\noindent\destacado{Sobre los programas de la PUCP.} Se incorporan los temas y
lecturas de siete programas de la PUCP (Anexo E), salvo los de macroeconomía
internacional, finanzas y banca, y política macroeconómica peruana. Las líneas
\emph{Artículos (PUCP)} y \emph{Artículos} (lecturas que no proceden de los
programas) entran en el cómputo de horas; \emph{Consulta} y \emph{Consulta
(PUCP)}, no. Referencias completas en el Anexo F.

\medskip
\noindent\destacado{Sobre el texto suplementario.} \textcolor{colorK}{\textbf{K}}
(Azzimonti, Krusell, McKay y Mukoyama) es un manual de nivel doctoral y no forma
parte del recorrido obligatorio. Aparece únicamente donde profundiza un tema que
\emph{ya} está en el plan, y siempre como lectura posterior a la de los textos
base: primero se cierra el tema con De Gregorio y Heijdra, y solo entonces se
abre el suplemento. Se han omitido deliberadamente los capítulos de \textcolor{colorK}{\textbf{K}}
que introducen materia nueva respecto de este sílabo; véase el Anexo C.

\medskip
\noindent\destacado{Notación.} \DG{n} remite al capítulo $n$ de De Gregorio, y
\HJ{n}, \RO{n}, \CSV{n}, \AC{n}, \WA{n}, \GA{n} y \KR{n}, al de Heijdra, Romer,
Campante et al., Acemoglu, Walsh, Galí y el suplemento; \GM{n}, al libro de
consulta de Galindo y Montecinos (sin horas). Si solo entra parte de un capítulo
se indican secciones: \RO{8.6} es la sección 8.6.




%------------------------------------------------------------------------------
\newpage
\section*{Cronograma por semanas}
\addcontentsline{toc}{section}{Cronograma por semanas}

Calendario del plan con \textbf{suplemento integrado}, excluido el material
introductorio: 1443 horas distribuidas en \textbf{85 semanas de contenido}, que
equivalen a unos \textbf{22 meses de calendario} una vez descontadas las semanas
perdidas por feriados, viajes e imprevistos. Las horas de los bloques revisados
(1, 4, 5, 6 y 12 a 15) se reestimaron en proporción a las lecturas con horas que
conservan o suman.

\medskip
\noindent Todas las semanas son de materia nueva: no hay semanas de repaso ni
colchones intermedios. La columna de horas es la que gobierna el reparto; el
número de semanas es su traducción aproximada a un ritmo de diecisiete horas
semanales, de modo que algunos bloques corren ligeramente por encima o por
debajo de esa cifra.

\medskip
{\filasalternas
\noindent\begin{tabularx}{\textwidth}{lXcc}
\toprule
\rowcolor{tonoOscuro}\cab{Semanas} & \cab{Contenido} & \cab{Sem.} & \cab{Horas} \\
\midrule
1--3 & Bloque 1 --- El corto plazo tradicional & 3 & 55 \\
4--10 & Bloque 2 --- Microfundamentos & 7 & 111 \\
11--13 & Bloque 3 --- Equilibrio real y economía abierta & 3 & 54 \\
\rowcolor{tonoHielo}\textbf{14--20} & \textbf{Bloque 4 --- Crecimiento exógeno y ahorro óptimo} & \textbf{7} & \textbf{122} \\
21--24 & Bloque 5 --- Crecimiento endógeno y cambio tecnológico & 4 & 71 \\
25--28 & Bloque 6 --- Desarrollo, instituciones y economía política & 4 & 66 \\
29--32 & Bloque 7 --- Generaciones solapadas & 4 & 71 \\
\rowcolor{tonoHielo}\textbf{33--41} & \textbf{Bloque 8 --- Dinero, inflación y política monetaria} & \textbf{9} & \textbf{160} \\
42--46 & Bloque 9 --- Expectativas racionales y credibilidad & 5 & 86 \\
47--51 & Bloque 10 --- El mercado del trabajo & 5 & 78 \\
52--55 & Bloque 11 --- Rigideces nominales & 4 & 71 \\
56--60 & Bloque 14 --- Métodos recursivos y cuantitativos (temas 14.1 a 14.7) & 5 & 81 \\
\rowcolor{tonoHielo}\textbf{61--74} & \textbf{Bloque 12 --- Equilibrio general dinámico y estocástico (con el tema 14.8)} & \textbf{14} & \textbf{241} \\
75--79 & Bloque 13 --- Política monetaria neokeynesiana & 5 & 77 \\
80--85 & Bloque 15 --- Fricciones y heterogeneidad en el ciclo & 6 & 99 \\
\midrule
\rowcolor{tonoHielo}\multicolumn{2}{l}{\textbf{Total}} & \textbf{85} & \textbf{1443} \\
\bottomrule
\end{tabularx}
}

\medskip
\noindent\destacado{Ritmo supuesto.} Unas diecisiete horas por semana, todas de
contenido nuevo: por ejemplo, seis sesiones de dos horas de lunes a sábado y una
sesión larga de cinco horas.

\medskip
\noindent\destacado{El orden de lectura.} El Bloque 14 se estudia antes que el 12:
primero las herramientas (programación dinámica, equilibrio recursivo,
soluciones locales y cómputo) y después los modelos DSGE, que así se pueden
resolver y simular mientras se estudian. El tema 14.8 va dentro del Bloque 12,
después del 12.7.

\medskip
\noindent\destacado{Los tramos largos.} Los Bloques 12, 8 y 4 son los más largos:
suman 30 de las 85 semanas. Además, el crecimiento (Bloques 4 a 6, semanas
14--28) y el tramo final (Bloque 14 y Bloques 12, 13 y 15, semanas 56--85)
forman dos recorridos seguidos de alta exigencia. Son los puntos donde el avance suele atascarse, y al no haber
colchón posterior cualquier retraso acumulado se arrastra hasta el final del
plan.

\medskip
\noindent\destacado{Hitos de control.} Conviene verificar el avance al terminar las
semanas 13, 28, 55 y 74. Si en alguno de esos cortes se acumula más de una
semana de retraso, el ajuste correcto es alargar el calendario y no comprimir el
bloque siguiente.

\medskip
\noindent\destacado{Consolidación.} Sin sesiones dedicadas al repaso, el
mantenimiento de lo ya estudiado pasa a depender de los ejercicios del manual de
Heijdra, que retoman material de capítulos anteriores. Conviene resolverlos a
medida que avanza cada bloque y no dejarlos para el final.

%------------------------------------------------------------------------------
\newpage
\bloque{Bloque 1 --- El corto plazo tradicional --- 4 semanas}
{De Gregorio y Heijdra tratan este material como conocido, así que se despacha
rápido: fija el lenguaje común y, sobre todo, aplica las herramientas dinámicas
—diagramas de fase, sendero de silla— que se usarán en todos los bloques
posteriores. Su matemática está en los Bloques 9 y 10 del sílabo de economía
matemática.}

\tema{1.1 El modelo IS-LM de economía cerrada}
\lecturas{\DG{19} \quad \HJ{1}}
\consultalib{\CSV{15.1} \quad Jiménez (2007), caps.~4 y 7}

\tema{1.2 Economía abierta de corto plazo: el modelo de Mundell--Fleming}
\lecturas{\DG{20} \quad \HJ{2}}
\nota{Regímenes de tipo de cambio fijo y flexible; grado de movilidad de capitales.}

\tema{1.3 Sobrerreacción del tipo de cambio: el modelo de Dornbusch}
\lecturas{\HJ{2}}

\tema{1.4 Oferta agregada, curva de Phillips y tasa natural de desempleo}
\lecturas{\DG{21} \quad \HJ{3} \quad \RO{6.1--6.3}}
\nota{Romer añade la pregunta clave del bloque: si la rigidez está en precios o
en salarios, y qué implica para el comportamiento cíclico del salario real.}

\tema{1.5 Equilibrio AD-AS y políticas macroeconómicas}
\lecturas{\DG{22} \quad \HJ{3}}
\consultalib{\CSV{15.1} \quad \RO{6.4} \quad Jiménez (2007), cap.~10}

\tema{1.6 Previsión perfecta, diagramas de fase y sendero de silla}
\lecturas{\HJ{4}}
\nota{Punto de inflexión técnico del curso. Política anunciada frente a política
no anticipada; efectos de anuncio. Los diagramas de fase y la linealización se
estudian en los temas 10.2 y 10.4 del sílabo de economía matemática.}

%------------------------------------------------------------------------------
\bloque{Bloque 2 --- Microfundamentos: consumo, inversión y gobierno --- 7 semanas}
{Se abandonan las funciones de comportamiento \emph{ad hoc} y se derivan las
decisiones a partir de la optimización de hogares y empresas.}

\tema{2.1 Consumo intertemporal: ciclo de vida e ingreso permanente}
\lecturas{\DG{3} \quad \CSV{11} \quad \RO{8.1} \quad \RO{8.4}}
\supl{\KR{11}}
\consultalib{Jiménez (2007), cap.~3 \quad Prácticas UP (2024), secc.~2}
\nota{Restricciones de liquidez y evidencia empírica. La versión estocástica
—paseo aleatorio de Hall— se ve en el Bloque 12 con \HJ{17}.}

\tema{2.2 Inversión, costo de uso del capital y costos de ajuste}
\lecturas{\DG{4} \quad \CSV{13.1} \quad \HJ{4} \quad \RO{9.1} \quad \RO{9.7--9.8}}
\consultalib{Prácticas UP (2024), secc.~2}
\nota{Romer y Campante et al. añaden la inversión bajo incertidumbre: el valor de
opción de esperar (crítica de Pindyck) y los costos de ajuste fijos o con quiebre.}

\tema{2.3 La \emph{q} de Tobin y la dinámica de la inversión}
\lecturas{\HJ{4} \quad \DG{4} \quad \CSV{13.2} \quad \RO{9.2--9.6} \quad \AC{7.8}}
\articulos{Hayashi (1982)}
\nota{Salto de dificultad notable respecto de \DG{4}: es el primer modelo
dinámico plenamente resuelto del curso.}

\tema{2.4 El gobierno: restricción presupuestaria intertemporal}
\lecturas{\DG{5} \quad \CSV{17.1} \quad \HJ{6} \quad \RO{13.1}}
\supl{\KR{15}}

\tema{2.5 Equivalencia ricardiana y sus límites}
\lecturas{\HJ{6} \quad \DG{5} \quad \CSV{17.2} \quad \RO{13.2}}
\consultalib{Prácticas UP (2024), secc.~2 \quad Ljungqvist y Sargent (2018), cap.~10}
\nota{La razón más importante por la que falla —horizontes finitos— se entiende
recién en el Bloque 7, con \HJ{15}.}

\tema{2.6 Sostenibilidad de la deuda pública y señoreaje}
\lecturas{\HJ{6} \quad \DG{5} \quad \DG{16} \quad \CSV{17.7--17.8} \quad \RO{13.8}}
\supl{\KR{15}}
\nota{Campante et al. tratan los problemas de medición del déficit: inflación,
venta de activos, pasivos contingentes y enfoque de balance.}

\tema{2.7 Suavizamiento de impuestos}
\lecturas{\CSV{18.1} \quad \RO{13.3}}
\consultalib{Prácticas UP (2024), secc.~4}
\nota{El modelo de Barro: por qué conviene financiar con deuda los gastos
transitorios. La imposición óptima del capital en el modelo neoclásico
(\CSV{18.3}) se estudia en el tema 4.8.}

\tema{2.8 Economía política de los déficits y reglas fiscales}
\lecturas{\CSV{18.2} \quad \RO{13.4--13.7}}
\nota{Acumulación estratégica de deuda, estabilización demorada y evidencia sobre
política y déficits en países industrializados; reglas e instituciones
fiscales.}

%------------------------------------------------------------------------------
\bloque{Bloque 3 --- Equilibrio real y economía abierta intertemporal --- 3 semanas}
{Bloque casi exclusivo de De Gregorio y su mayor aporte diferencial. Heijdra no
desarrolla el tipo de cambio real con esta profundidad.}

\tema{3.1 La economía cerrada de pleno empleo}
\lecturas{\DG{6}}

\tema{3.2 La cuenta corriente como fenómeno intertemporal}
\lecturas{\DG{7}}
\supl{\KR{23}}

\tema{3.3 El tipo de cambio real: PPP, bienes transables y no transables}
\lecturas{\DG{8}}

\tema{3.4 Determinantes del tipo de cambio real y de la cuenta corriente}
\lecturas{\DG{9}}
\supl{\KR{23}}
\nota{Efecto Balassa--Samuelson, términos de intercambio, cuenta corriente de
equilibrio.}

\tema{3.5 Complemento formal: economía abierta pequeña con ahorro óptimo}
\lecturas{\HJ{13} (secciones de economía abierta)}
\nota{Se aprovecha mejor una vez estudiado el Bloque 4.}

%------------------------------------------------------------------------------
\bloque{Bloque 4 --- Crecimiento exógeno y ahorro óptimo --- 11 semanas}
{Núcleo del curso. Contiene el paso decisivo hacia la macroeconomía moderna: el
hogar que optimiza y la ecuación de Euler.
El crecimiento endógeno, el desarrollo y la economía política del crecimiento
continúan en los Bloques 5 y 6, con Acemoglu como texto principal.}

\tema{4.1 Hechos estilizados y contabilidad del crecimiento}
\lecturas{\DG{10} \quad \RO{1.1} \quad \AC{1} \quad \AC{3.1}}
\consultalib{\CSV{2.1--2.2}}
\articulos{Jones (2016)}
\consulta{Williamson (2010), caps.~7--8 \quad Duncan (2022)}
\nota{Hechos de Kaldor y diferencias de ingreso entre países; medición, debates y
mitos sobre el crecimiento (programa 1ECO74).}

\tema{4.2 El modelo de Solow--Swan: estado estacionario y regla dorada}
\lecturas{\DG{11} \quad \HJ{12} \quad \AC{2.1--2.4} \quad \AC{2.6--2.7}}
\consultalib{\CSV{2.3} \quad \CSV{2.5} \quad \RO{1.2--1.4} \quad \KR{3}}
\articulos{Solow (1956)}
\consulta{Barro y Sala-i-Martin (2003)}
\nota{Acemoglu lo desarrolla en tiempo discreto y continuo, con progreso técnico y
estática comparativa.}

\tema{4.3 Dinámica de transición, convergencia y velocidad de ajuste}
\lecturas{\DG{12} \quad \HJ{12} \quad \AC{2.5}}
\consultalib{\RO{1.5--1.6} \quad \KR{3}}
\consulta{Duncan y Fuentes (2006)}

\tema{4.4 Capital humano: la ampliación de Mankiw--Romer--Weil}
\lecturas{\DG{12} \quad \HJ{12} \quad \AC{3.3}}
\consultalib{\CSV{2.4} \quad \RO{4.1}}

\tema{4.5 Evidencia empírica sobre el crecimiento}
\lecturas{\DG{13} \quad \RO{1.7} \quad \RO{4.2} \quad \AC{3.2} \quad \AC{3.4--3.6}}
\consultalib{\CSV{7.1} \quad \KR{13} \quad \KR{8}}
\nota{Sin contraparte en Heijdra. Regresiones de convergencia y determinantes
del crecimiento.}
\nota{Acemoglu añade la calibración y la estimación de las diferencias de
productividad entre países.}

\tema{4.6 Fundamentos del crecimiento neoclásico: hogar representativo y teoremas del bienestar}
\lecturas{\AC{5}}
\nota{Horizonte infinito, empresa representativa, teoremas del bienestar y comercio
secuencial: la base del problema del planificador del tema 4.7.}

\tema{4.7 Ahorro óptimo: el modelo de Ramsey--Cass--Koopmans}
\lecturas{\DG{14} \quad \HJ{13} \quad \AC{8.1--8.7} \quad \AC{8.9--8.10}}
\consultalib{\CSV{3} \quad \RO{2.1--2.6} \quad \KR{4} \quad \KR{5} \quad Stokey, Lucas y Prescott (1989), secc.~5.4}
\nota{Ecuación de Euler, elasticidad de sustitución intertemporal, condición de
transversalidad. Supone el control óptimo del Bloque 12 del sílabo de economía
matemática: sin él, la ecuación de Euler se convierte en una fórmula
memorizada. \textbf{No avanzar sin dominar este tema.}}

\tema{4.8 Eficiencia dinámica y política fiscal en el modelo de Ramsey}
\lecturas{\HJ{13} \quad \CSV{18.3} \quad \AC{8.8}}
\consultalib{\RO{2.7} \quad \KR{6} \quad Ljungqvist y Sargent (2018), cap.~11 (secc.~11.1--11.6) y cap.~16 (secc.~16.1--16.7)}
\consulta{Mankiw y Weinzierl (2006)}
\nota{Efectos del gasto público e imposición óptima del capital en el modelo
neoclásico. Ljungqvist y Sargent resuelven el modelo con impuestos
distorsionantes (cap.~11) y el problema de Ramsey con el impuesto cero al
capital de Chamley y Judd (cap.~16), que enlaza con el Bloque 9 del sílabo de
economía pública.}

\tema{4.9 Crecimiento, recursos naturales y medio ambiente}
\lecturas{\RO{1.8}}
\supl{\KR{25}}
\consulta{NobelPrize.org (2018)}
\nota{Límites ambientales al crecimiento y vínculos entre crecimiento, cambio
tecnológico y clima.}

%------------------------------------------------------------------------------
\bloque{Bloque 5 --- Crecimiento endógeno y cambio tecnológico --- 5 semanas}
{Lleva el crecimiento endógeno más allá del modelo AK: capital humano, innovación,
variedades, calidad y sesgo del cambio técnico. Arranca con un panorama del
crecimiento endógeno (tema 5.1); en el resto, Acemoglu es la lectura
principal. Requiere el Bloque 4.}

\tema{5.1 Crecimiento endógeno: modelos AK, externalidades e I+D}
\lecturas{\DG{12} \quad \HJ{14} \quad \AC{11}}
\consultalib{\CSV{5.1--5.4} \quad \RO{3.1--3.3} \quad \KR{13} \quad Stokey, Lucas y Prescott (1989), secc.~5.6--5.8}
\nota{Romer, Lucas, variedad de productos; política pública y crecimiento.}
\nota{Es el panorama del bloque; los temas 5.3 a 5.6 desarrollan cada mecanismo.}

\tema{5.2 Capital humano y crecimiento}
\lecturas{\AC{10}}
\nota{Teorema de separación, inversión en escolaridad, modelo de Ben-Porath,
complementariedad entre capital y calificación, externalidades de capital humano
y el modelo de Nelson--Phelps.}

\tema{5.3 El modelado del cambio tecnológico}
\lecturas{\RO{3.4} \quad \AC{12}}
\nota{Concepciones de la tecnología, valor de la innovación y el agregador de
Dixit--Stiglitz aplicado a la innovación.}

\tema{5.4 Variedades crecientes: el modelo de Romer}
\lecturas{\CSV{6.1} \quad \RO{3.5} \quad \AC{13}}
\nota{Insumos de variedad creciente, derrames de conocimiento y crecimiento sin efectos
de escala.}

\tema{5.5 Crecimiento schumpeteriano y destrucción creativa}
\lecturas{\CSV{6.2--6.3} \quad \AC{14}}
\consultalib{\KR{22}}
\nota{Escaleras de calidad, innovación de incumbentes y entrantes, distancia a la
frontera tecnológica y competencia.}

\tema{5.6 Cambio tecnológico dirigido}
\lecturas{\AC{15}}
\nota{Por qué el cambio técnico favorece a ciertos factores; cambio técnico aumentador
de trabajo.}

\tema{5.7 Efectos de escala, población y el futuro del crecimiento}
\lecturas{\CSV{6.4} \quad \RO{3.6--3.8}}
\nota{Pruebas de series de tiempo del crecimiento endógeno; población y cambio
tecnológico en el muy largo plazo.}

%------------------------------------------------------------------------------
\bloque{Bloque 6 --- Desarrollo, instituciones y economía política del crecimiento --- 5 semanas}
{Pasa de los mecanismos del crecimiento a sus causas fundamentales: difusión de
tecnología, cambio estructural, trampas de pobreza e instituciones. Recorre las
partes de Acemoglu sobre desarrollo y economía política. Requiere el Bloque 5.
Los temas 6.6 y 6.7, de instituciones y economía política, son opcionales: no
entran en el cómputo de horas y pueden leerse junto con el sílabo de economía
pública.}

\tema{6.1 Causas próximas y fundamentales del crecimiento}
\lecturas{\CSV{7.2} \quad \RO{4.3--4.5} \quad \AC{4}}
\consulta{Ogilvie y Carus (2014) \quad Bloom, Kuhn y Prettner (2022)}
\nota{Suerte, geografía, cultura e instituciones; enfermedad y desarrollo.}

\tema{6.2 Difusión de tecnología y diferencias de productividad}
\lecturas{\RO{4.6} \quad \AC{18}}
\articulos{Syverson (2011)}
\consultalib{\AC{19.1--19.7}}
\nota{Tecnologías apropiadas e inapropiadas, e instituciones de contratación en la
adopción tecnológica.}
\nota{\AC{19} extiende el tema a la economía abierta: por qué el capital no
fluye de los países ricos a los pobres, comercio y especialización, ciclo del
producto y aprendizaje en la práctica. Es consulta, sin horas; el comercio se
estudia con horas en el sílabo de economía internacional.}

\tema{6.3 Cambio estructural y crecimiento no balanceado}
\lecturas{\AC{20}}
\nota{Fuentes de demanda y de oferta del crecimiento no balanceado; productividad
agrícola e industrialización.}

\tema{6.4 Transformaciones estructurales, trampas de pobreza y fallas de mercado}
\lecturas{\CSV{5.5} \quad \AC{17.6} \quad \AC{21.1--21.6}}
\articulos{Balboni et al. (2021)}
\nota{Desarrollo financiero, transición demográfica, economía dual, gran empuje y
desigualdad con mercados de crédito imperfectos.}

\tema{6.5 Teoría unificada: de Malthus al crecimiento moderno}
\lecturas{\CSV{10} \quad \AC{21.7} \quad \AC{24.2}}
\consulta{P. Romer (1996)}
\nota{Régimen malthusiano, transición demográfica y despegue del crecimiento
sostenido.}

\tema{6.6 Instituciones, conflicto distributivo y crecimiento (opcional)}
\consultalib{\AC{22}}
\nota{Conflicto distributivo, instituciones económicas ineficientes, votante mediano y
provisión de bienes públicos con Estados débiles o fuertes.}

\tema{6.7 Instituciones políticas y crecimiento (opcional)}
\consultalib{\AC{23} \quad \AC{24.1} \quad \AC{24.3} \quad \AC{Apéndice C}}
\nota{Regímenes políticos, cambio político endógeno y dinámica de las instituciones.
El apéndice C repasa los juegos dinámicos que usa el capítulo.}

%------------------------------------------------------------------------------
\bloque{Bloque 7 --- Generaciones solapadas --- 4 semanas}
{Sin equivalente en De Gregorio. Heijdra lleva el peso formal; Campante et al.
y Romer ofrecen versiones compactas del modelo de Diamond y de las reformas de
pensiones, y Acemoglu (cap.~9), el tratamiento más completo.}

\tema{7.1 El modelo de Diamond en tiempo discreto}
\lecturas{\CSV{8.1} \quad \HJ{16} \quad \RO{2.8--2.10} \quad \AC{9.1--9.3}}

\tema{7.2 Ineficiencia dinámica y sobreacumulación de capital}
\lecturas{\CSV{8.2} \quad \HJ{16} \quad \RO{2.11} \quad \AC{9.4}}
\supl{\KR{6}}

\tema{7.3 Seguridad social: reparto frente a capitalización}
\lecturas{\CSV{9.1} \quad \HJ{16} \quad \AC{9.5}}

\tema{7.4 El modelo de Blanchard--Yaari en tiempo continuo}
\lecturas{\CSV{8.3} \quad \HJ{15} \quad \AC{9.6--9.8}}
\nota{Probabilidad de muerte constante, agregación, y el motivo por el cual falla
la equivalencia ricardiana.}
\nota{Acemoglu añade el altruismo imperfecto como motivo de herencia.}

\tema{7.5 Envejecimiento poblacional, pensiones y deuda pública}
\lecturas{\CSV{9.2} \quad \HJ{15} \quad \HJ{16} \quad \RO{2.12}}
\nota{Campante et al. analizan cómo financiar la transición desde un sistema de
reparto: impuestos a los jóvenes o emisión de deuda.}

\tema{7.6 Burbujas racionales y la deuda pública como burbuja}
\lecturas{\CSV{22.3}}
\nota{Solo hay burbujas sostenibles si la economía es dinámicamente ineficiente
(tema 7.2). Implicaciones para la política fiscal y financiera.}

%------------------------------------------------------------------------------
\bloque{Bloque 8 --- Dinero, inflación y política monetaria --- 9 semanas}
{De Gregorio cubre la política monetaria aplicada, y Heijdra, la integración del
dinero en el equilibrio general.
Walsh es la referencia formal de todo el bloque; Galí aporta el modelo
monetario clásico, y Campante et al. y Romer, versiones compactas.}

\tema{8.1 La naturaleza del dinero y las nuevas formas monetarias}
\lecturas{\WA{3.4}}
\articulos{Kiyotaki y Wright (1989) \quad Benigno (2025), cap.~3}
\consultalib{Ljungqvist y Sargent (2018), secc.~29.8}
\consulta{Lewis y Mizen (2000), caps.~1--3 \quad Mishkin (2021), cap.~3 \quad Kiyotaki y Moore (2002) \quad McCallum (2015) \quad Böhme et al. (2015) \quad Velde (2013) \quad BIS (2018) \quad BIS (2025)}
\nota{Funciones del dinero, dinero fiduciario y crediticio y origen de la banca
central; dinero electrónico, criptomonedas y monedas digitales de banco central.
El modelo de búsqueda de Kiyotaki y Wright explica por qué surge un medio de
cambio.}

\tema{8.2 Dinero, precios y producto: la evidencia}
\lecturas{\RO{5.9} \quad \WA{1}}
\consultalib{Sørensen (1996) \quad Stock y Watson (2001)}
\nota{Correlaciones de largo y corto plazo, evidencia de Friedman y Schwartz,
causalidad de Granger y vectores autorregresivos.}

\tema{8.3 Teoría cuantitativa, neutralidad y demanda por dinero}
\lecturas{\DG{15} \quad \HJ{10} \quad \WA{3.1--3.2}}
\supl{\KR{17}}
\articulos{Baumol (1952) \quad Tobin (1956)}
\consulta{Tobin (1958) \quad Mishkin (2021), cap.~19 \quad Quispe (1998) \quad León (1999)}
\consultalib{Jiménez (2007), cap.~5}
\nota{Demanda por transacciones de Baumol--Tobin, modelos de tiempo de compra y
estabilidad empírica de la demanda de dinero.}

\tema{8.4 Oferta de dinero, señoreaje, inflación e hiperinflaciones}
\lecturas{\DG{16} \quad \CSV{19.3} \quad \HJ{10} \quad \RO{12.10} \quad \WA{4.1--4.4}}
\articulos{Dornbusch (1992)}
\consultalib{Ljungqvist y Sargent (2018), cap.~27, secc.~27.1--27.3, que sirve también para los temas 8.6 y 8.7}
\consulta{Lewis y Mizen (2000), cap.~7 \quad Woodford (2003), cap.~6}
\consultalib{Jiménez (2007), cap.~6}
\nota{Restricción presupuestaria consolidada, señoreaje de equilibrio, modelo de
Cagan, hiperinflación racional y aritmética monetarista desagradable.}

\tema{8.5 El dinero en equilibrio general: MIU, \emph{cash-in-advance} y regla de Friedman}
\lecturas{\CSV{19.1--19.2} \quad \HJ{10} \quad \WA{2.1--2.2} \quad \WA{2.4--2.8} \quad \WA{3.3} \quad \WA{3.5--3.7} \quad \GA{2}}
\supl{\KR{17}}
\articulos{Benigno (2025), cap.~6}
\consultalib{Ljungqvist y Sargent (2018), cap.~9, secc.~9.1--9.5}
\consulta{McCallum y Nelson (2010) \quad Lucas (2013) \quad Woodford (2003), caps.~2--3}
\nota{Incluye generaciones solapadas con dinero; conviene haber cursado el Bloque 7.}
\nota{Walsh resuelve y calibra el modelo MIU; Galí presenta el modelo monetario
clásico que sirve de base al neokeynesiano del Bloque 12.}

\tema{8.6 Costos de la inflación y cantidad óptima de dinero}
\lecturas{\CSV{19.4} \quad \WA{2.3} \quad \WA{4.6--4.8}}
\articulos{Lucas (1994)}
\consulta{Lucas (1996)}
\nota{Costo de bienestar de la inflación, señoreaje óptimo y la regla de Friedman
revisitada; inflación y competencia (modelo de Tommasi).}

\tema{8.7 La teoría fiscal del nivel de precios}
\lecturas{\CSV{22.2} \quad \WA{4.5}}
\nota{Políticas fiscales ricardianas y no ricardianas: cuándo la restricción
presupuestaria del gobierno determina el nivel de precios.}

\tema{8.8 Política monetaria, tasas de interés y mercados financieros}
\lecturas{\DG{17} \quad \RO{12.1--12.2} \quad \WA{10.1} \quad \WA{10.3} \quad \WA{10.6--10.8}}
\articulos{Bernanke y Gertler (1995)}
\consulta{Mishkin (1995) \quad Benigno (2025), cap.~8 \quad Mishkin (2021), caps.~4--7}
\consultalib{Jiménez (2007), cap.~7}
\nota{Sin contraparte en Heijdra. Mecanismos de transmisión y estructura temporal
de tasas.}
\nota{Hipótesis de expectativas de la estructura temporal y canal del crédito
(préstamos bancarios y hoja de balance).}

\tema{8.9 Procedimientos operativos de la política monetaria}
\lecturas{\WA{11}}
\articulos{Poole (1970)}
\consulta{Goodhart (1994) \quad Bindseil (2014), caps.~1--3 \quad Lewis y Mizen (2000), cap.~16 \quad Mishkin (2021), caps.~14--15}
\consultalib{Poole (1968)}
\nota{Elección de instrumento según Poole, metas intermedias, mercado de reservas y
sistema de corredor.}
\nota{Son dos trabajos distintos de Poole: el de 1970 trata la elección del
instrumento; el de 1968, la demanda de reservas de un banco bajo incertidumbre,
que es el modelo del mercado interbancario y del corredor de tasas que usa el
curso de Política Monetaria del BCRP.}

%------------------------------------------------------------------------------
\bloque{Bloque 9 --- Expectativas racionales y credibilidad --- 5 semanas}
{Los tres capítulos base de este bloque (\HJ{5}, \HJ{9} y \DG{25}) se leen muy
bien de forma consecutiva. Walsh aporta el tratamiento formal de las rigideces
informativas y del sesgo inflacionario.}

\tema{9.1 Expectativas racionales y métodos de solución}
\lecturas{\HJ{5}}

\tema{9.2 Proposición de inefectividad de la política y crítica de Lucas}
\lecturas{\HJ{5} \quad \CSV{15.2} \quad \RO{6.9}}
\articulos{Lucas (1976)}
\consultalib{Jiménez (2007), cap.~8}
\nota{El modelo de información imperfecta de Lucas (las «islas») da la curva de
oferta sobre la que descansa la proposición de inefectividad.}

\tema{9.3 Rigideces informativas y de cartera}
\lecturas{\WA{5} \quad \WA{6.3.3}}
\articulos{Mankiw y Reis (2010)}
\consulta{Sims (2010)}
\nota{Información imperfecta e información rígida, aprendizaje, y participación
limitada con efectos de liquidez.}

\tema{9.4 Inconsistencia intertemporal y sesgo inflacionario}
\lecturas{\DG{25} \quad \CSV{20.1} \quad \HJ{9} \quad \RO{12.8} \quad \WA{7.1--7.2}}
\articulos{Barro y Gordon (1983)}
\consultalib{Kydland y Prescott (1977)}
\nota{Kydland y Prescott (1977) plantean el problema general: el plan óptimo
anunciado deja de ser óptimo una vez que el público fijó sus expectativas.
Barro y Gordon (1983) lo aplican a la política monetaria.}

\tema{9.5 Reputación, delegación e independencia del banco central}
\lecturas{\HJ{9} \quad \DG{25} \quad \CSV{20.2} \quad \RO{12.9} \quad \WA{7.3--7.6}}
\articulos{Rogoff (1985) \quad Walsh (1995) \quad Bernanke y Mishkin (1997)}
\consultalib{Ljungqvist y Sargent (2018), cap.~24, secc.~24.1--24.7}
\consulta{Reis (2013) \quad Mishkin (2000) \quad King (1999) \quad Mishkin (2021), caps.~16 y 23 \quad Lewis y Mizen (2000), caps.~10 y 15}
\nota{Banquero central conservador, contratos óptimos, metas de inflación.}

%------------------------------------------------------------------------------
\bloque{Bloque 10 --- El mercado del trabajo --- 5 semanas}
{El peso analítico recae en Heijdra y Romer. La extensión de estos modelos al
ciclo económico se estudia en el Bloque 15.}

\tema{10.1 El mercado del trabajo y el mercado del crédito en el ciclo}
\lecturas{\DG{24} \quad \CSV{16.1}}
\supl{\KR{19}}

\tema{10.2 Salarios de eficiencia}
\lecturas{\CSV{16.4} \quad \HJ{7} \quad \RO{11.1--11.2}}
\nota{Romer desarrolla el modelo de Shapiro--Stiglitz completo.}

\tema{10.3 Sindicatos, negociación colectiva y modelos \emph{insider--outsider}}
\lecturas{\CSV{16.5} \quad \HJ{7}}
\nota{Campante et al. incluyen la migración rural-urbana como fuente de desempleo.}

\tema{10.4 Búsqueda y emparejamiento: el modelo Diamond--Mortensen--Pissarides}
\lecturas{\CSV{16.2--16.3} \quad \HJ{8} \quad \RO{11.4--11.6}}
\supl{\KR{20}}
\articulos{Mortensen y Pissarides (1994)}
\consultalib{Ljungqvist y Sargent (2018), cap.~6, secc.~6.1--6.4 (modelo de McCall y \emph{lake model}) \quad Stokey, Lucas y Prescott (1989), secc.~10.7--10.8}
\consulta{Pissarides (2000) \quad Davis y Haltiwanger (1992)}
\nota{Función de emparejamiento, curva de Beveridge, negociación de Nash del
salario, creación y destrucción de empleos.}

\tema{10.5 Contratos implícitos y rigidez salarial}
\lecturas{\RO{11.3}}
\nota{Contratos entre empresas y trabajadores adversos al riesgo: por qué el
salario puede no responder al ciclo aunque el mercado no tenga fricciones de
búsqueda.}

%------------------------------------------------------------------------------
\bloque{Bloque 11 --- Rigideces nominales con microfundamentos --- 4 semanas}
{Heijdra y Romer llevan el peso del bloque, que es prerrequisito conceptual del
Bloque 12: sin competencia monopolística el modelo neokeynesiano no es
interpretable.}

\tema{11.1 Competencia monopolística: el agregador de Dixit--Stiglitz}
\lecturas{\CSV{15.3} \quad \HJ{11} \quad \RO{6.5} \quad \WA{6.2.3}}
\articulos{Blanchard y Kiyotaki (1987)}

\tema{11.2 Costos de menú y rigidez nominal de precios}
\lecturas{\HJ{11} \quad \RO{6.6} \quad \WA{6.2.5}}
\supl{\KR{18}}

\tema{11.3 Rigideces reales frente a rigideces nominales}
\lecturas{\CSV{15.3} \quad \HJ{11} \quad \RO{6.7} \quad \WA{6.1} \quad \WA{6.2.1}}

\tema{11.4 Fallas de coordinación y equilibrios múltiples}
\lecturas{\RO{6.8}}
\nota{Complementariedades estratégicas y teorías reales no walrasianas de las
fluctuaciones.}

\tema{11.5 Evidencia microeconómica sobre la fijación de precios}
\lecturas{\RO{7.5--7.6} \quad \WA{6.2.6} \quad \WA{6.3.1}}
\articulos{Nakamura y Steinsson (2008) \quad Klenow y Malin (2010)}
\consulta{Nakamura y Steinsson (2013) \quad Midrigan (2011) \quad Gorodnichenko y
Talavera (2017) \quad Cavallo (2018) \quad Cavallo, Lippi y Miyahara (2024)
\quad Bueno (2024)}
\nota{Frecuencia y tamaño de los cambios de precios; fijación dependiente del
estado frente a dependiente del tiempo. Bueno (2024) usa precios de
supermercados del Perú.}

\tema{11.6 Multiplicadores del gasto público en equilibrio general}
\lecturas{\HJ{11} \quad \CSV{17.3--17.4}}
\consulta{Galí, López-Salido y Vallés (2007)}
\nota{La referencia neoclásica —el gasto público en el modelo de Ramsey— está en el
tema 4.8. Galí et al. (2007) se aprovecha mejor después del tema 12.10.}

%------------------------------------------------------------------------------
\bloque{Bloque 12 --- Equilibrio general dinámico y estocástico (DSGE) --- 13 semanas}
{Núcleo cuantitativo del recorrido. Requiere haber cerrado los Bloques 4, 9 y 11
y los temas 14.1 a 14.7, que se estudian antes; el tema 14.8 (calibración y
simulación) se intercala en este bloque, después del 12.7. Los Bloques 13 y 15
lo prolongan. Es el bloque donde más pesan los programas de la PUCP, que lo
tratan como el contenido central de un primer año de doctorado.}

\tema{12.1 Hechos estilizados del ciclo económico}
\lecturas{\DG{18} \quad \RO{5.1}}
\articulos{Cooley y Prescott (1995)}
\consulta{Mitchell (1913)}
\consultalib{\GM{1.2}}
\nota{Comovimiento de producto, consumo, inversión y horas; medición y filtrado
del componente cíclico.}

\tema{12.2 Decisiones bajo incertidumbre: programación dinámica y ecuación de Bellman}
\lecturas{\HJ{17}}
\supl{\KR{7}}
\nota{Capítulo técnico y árido, pero sin él los dos últimos capítulos resultan
ilegibles.}
\nota{La teoría formal de la programación dinámica ya se estudió en los temas
14.1 y 14.2; aquí se aplica a las decisiones de los hogares.}

\tema{12.3 Consumo bajo incertidumbre y reparto competitivo de riesgos}
\lecturas{\CSV{12.1} \quad \HJ{17} \quad \RO{8.2--8.3}}
\supl{\KR{7} \quad \KR{11}}
\nota{Paseo aleatorio de Hall; cierra el tema abierto en 2.1.}

\tema{12.4 Más allá del ingreso permanente: ahorro precautorio y liquidez}
\lecturas{\CSV{12.2} \quad \RO{8.6--8.7}}
\articulos{Carroll (1997)}
\nota{\RO{8.7} resuelve el ahorro precautorio con programación dinámica: primera
aplicación directa del tema 12.2. Ahorro de amortiguación (\emph{buffer stock}) y
sesgo presente.}

\tema{12.5 El modelo de ciclo económico real (RBC)}
\lecturas{\DG{23} \quad \CSV{14.1} \quad \HJ{18} \quad \RO{5.2--5.7} \quad \AC{17.1--17.3}}
\supl{\KR{14} \quad \KR{12}}
\articulos{King, Plosser y Rebelo (1988)}
\consultalib{Stokey, Lucas y Prescott (1989), secc.~10.1}
\consulta{Kydland y Prescott (1982) \quad Long y Plosser (1983) \quad King y
Rebelo (1999) \quad Urrutia (1996)}
\consultalib{\GM{1.3--1.4} \quad \GM{5}}
\nota{De Gregorio da la intuición; Heijdra da la resolución completa, la
calibración y las funciones impulso-respuesta.}
\nota{Contexto histórico: de la crítica de Lucas (tema 9.2) a Kydland y Prescott
(1982). Senda de crecimiento balanceado y dinámica de transición.}
\nota{Acemoglu presenta el modelo de Brock--Mirman, la base estocástica del RBC.}

\tema{12.6 Extensiones del RBC: oferta de trabajo, costos de ajuste y utilización}
\lecturas{\CSV{14.1}}
\articulos{Hansen (1985)}
\consultalib{Ljungqvist y Sargent (2018), secc.~29.5--29.6 (loterías de empleo)}
\consulta{Rogerson (1988) \quad Greenwood, Hercowitz y Huffman (1988) \quad
Benhabib, Rogerson y Wright (1991)}
\consultalib{\GM{6} \quad Kim y Loungani (1992) \quad Rotemberg y Woodford (1996)
\quad Finn (2000) \quad Kilian (2008)}
\nota{Trabajo indivisible y loterías; hábitos en el consumo; costos de ajuste a la
inversión; utilización variable del capital; producción doméstica.}
\nota{RBC con energía: la energía entra como insumo de la producción y su precio,
como un choque de oferta persistente. Kim y Loungani (1992) lo incorporan al
RBC estándar y muestran que su efecto sobre el producto es modesto; Rotemberg y
Woodford (1996) lo amplifican con competencia imperfecta y Finn (2000) con
utilización variable del capital, sin abandonar la competencia perfecta. Kilian
(2008) resume la evidencia. Es el modelo que resuelve el Tema 3 del curso de
Macroeconomía del BCRP (estado estacionario, log-linealización y comparación con
un choque de productividad).}

\tema{12.7 Evaluación empírica de los modelos de ciclo real}
\lecturas{\HJ{18} \quad \DG{23} \quad \CSV{14.2--14.3} \quad \RO{5.8} \quad \RO{5.10}}
\supl{\KR{8} \quad \KR{14}}
\articulos{Prescott (1986)}
\consulta{Galí (1999) \quad Prescott (2016)}
\nota{Galí (1999) usa la respuesta de las horas a los choques tecnológicos como
prueba del modelo. Al cerrar este tema se estudia el 14.8: calibración y
simulación del RBC en la computadora.}

\tema{12.8 Contabilidad del ciclo económico: las cuñas}
\articulos{Chari, Kehoe y McGrattan (2007)}
\nota{Cuñas de eficiencia, de trabajo, de inversión y de gasto como diagnóstico
de qué fricción le falta al modelo. Se retoma en 15.5 con la cuña laboral.}

\tema{12.9 Ajuste escalonado de precios: Fischer, Taylor e inercia inflacionaria}
\lecturas{\RO{7.1--7.3} \quad \RO{7.7} \quad \WA{6.2.2} \quad \WA{6.2.4}}
\articulos{Fischer (1977)}
\nota{Precios predeterminados (Fischer) frente a precios fijos (Taylor); modelos
que generan la inercia que la curva de Phillips neokeynesiana pura no produce.}

\tema{12.10 El modelo DSGE neokeynesiano}
\lecturas{\CSV{15.4} \quad \HJ{19} \quad \RO{7.4} \quad \RO{7.8} \quad \WA{6.3.2} \quad \WA{8.1--8.3} \quad \GA{1} \quad \GA{3}}
\supl{\KR{18}}
\articulos{Calvo (1983)}
\consulta{Goodfriend y King (1997)}
\nota{Fijación de precios a la Calvo, curva de Phillips neokeynesiana, regla de
Taylor, modelo canónico de tres ecuaciones y política monetaria óptima.}
\nota{La política monetaria óptima en este modelo es el tema del Bloque 13. Su
simulación completa el tema 14.8.}

\tema{12.11 Modelos DSGE de mediana escala y su estimación}
\lecturas{\RO{7.10}}
\articulos{Smets y Wouters (2007)}
\nota{Hábitos, indexación, rigidez de salarios y costos de ajuste; estimación
bayesiana e identificación de choques.}

%------------------------------------------------------------------------------
\bloque{Bloque 13 --- Política monetaria en el modelo neokeynesiano --- 4 semanas}
{Desarrolla la política monetaria óptima en el modelo
neokeynesiano del tema 12.10. Es el núcleo de Galí y de los capítulos de
política de Walsh y Romer. Requiere los Bloques 8, 9 y 12.}

\tema{13.1 Diseño de la política monetaria en el modelo neokeynesiano básico}
\lecturas{\RO{12.3--12.5} \quad \WA{8.4.1--8.4.2} \quad \WA{8.5--8.7} \quad \GA{4}}
\articulos{Clarida, Galí y Gertler (1999)}
\consulta{Walsh (2002) \quad Woodford (2003)}
\consultalib{\GA{7} \quad \WA{9.1--9.5}}
\nota{La asignación eficiente como referencia, la función de pérdida derivada del
bienestar y las disyuntivas entre inflación y brecha del producto.}
\nota{\GA{7} y \WA{9} llevan el modelo a la economía abierta: la economía
pequeña de Galí--Monacelli, el modelo de dos países de Obstfeld--Rogoff y la
coordinación de políticas. Son consulta, sin horas; \GA{7} se lee con horas
en el tema 8.4 del sílabo de economía internacional.}

\tema{13.2 Reglas de tasa de interés, determinación del nivel de precios y principio de Taylor}
\lecturas{\RO{12.6} \quad \WA{10.2.1--10.2.2}}
\articulos{Taylor (1993) \quad Castillo-Martinez y Reis (2026)}
\consulta{Fernández-Villaverde y Sanches (2024)}
\nota{Reglas exógenas y de retroalimentación, unicidad del equilibrio, cláusulas de
escape y anclas; control de la inflación con tasas, con dinero y bajo dominancia
fiscal (programa 1ECO90).}

\tema{13.3 Discreción frente a compromiso}
\lecturas{\WA{8.4.3--8.4.7} \quad \GA{5}}
\nota{Ganancias del compromiso, persistencia endógena, regímenes de metas frente a
reglas de instrumento e incertidumbre de modelo.}

\tema{13.4 Salarios y precios rígidos}
\lecturas{\GA{6}}
\nota{Con dos fricciones nominales ya no basta con estabilizar la inflación de
precios.}

\tema{13.5 La cota inferior cero, la trampa de liquidez y el estancamiento secular}
\lecturas{\CSV{21.1} \quad \CSV{22.1} \quad \RO{12.7} \quad \WA{10.2.3}}
\articulos{Eggertsson y Woodford (2004)}
\consulta{Werning (2011) \quad Correia et al. (2013)}
\nota{Política monetaria y fiscal óptima cuando la tasa nominal no puede bajar más.}

\tema{13.6 \emph{Forward guidance} y política monetaria no convencional}
\lecturas{\CSV{21.2--21.4} \quad \RO{7.9} \quad \GA{8}}
\articulos{McKay, Nakamura y Steinsson (2016)}
\nota{El puzzle de \emph{forward guidance}, la hoja de balance del banco central y la
flexibilización cuantitativa y crediticia. Galí cierra con un balance del enfoque
neokeynesiano.}

\tema{13.7 Metas de inflación en una economía parcialmente dolarizada}
\articulosB{Céspedes, Chang y Velasco (2004) \quad Levy Yeyati (2006)}
\consultalib{Armas, Ize y Levy Yeyati (2006)}
\nota{La dolarización financiera como restricción de la política monetaria:
descalce de monedas, efectos de hoja de balance de una depreciación y por qué el
banco central interviene en el mercado cambiario y acumula reservas como segundo
instrumento junto a la tasa de referencia, que es la situación del BCRP.
Céspedes et al. muestran cuándo conviene el tipo de cambio flexible aun con
hojas de balance dolarizadas. El estudio con horas de la dolarización y las
crisis está en el tema 9.6 del sílabo de economía internacional; aquí se ve
desde la política monetaria.}

%------------------------------------------------------------------------------
\bloque{Bloque 14 --- Métodos recursivos y macroeconomía cuantitativa --- 7 semanas}
{Formaliza las herramientas de la macroeconomía dinámica y las lleva a la
computadora \emph{antes} de los modelos DSGE: los temas 14.1 a 14.7 se estudian
después del Bloque 11 y antes del 12, y el 14.8, dentro del Bloque 12, después
del tema 12.7. Es el núcleo común de los cursos avanzados de macroeconomía
dinámica de la PUCP (MAT616, 1ECO90 y ECO324). Requiere los Bloques 4 y 9 y,
del sílabo de economía matemática, el Bloque 13 (programación dinámica). El
peso recae en Stokey, Lucas y Prescott y en Ljungqvist y Sargent.}

\tema{14.1 Programación dinámica determinística: problema secuencial y ecuación funcional}
\articulos{Stokey, Lucas y Prescott (1989), caps.~1, 2, 4 y 5}
\consultalib{Stokey, Lucas y Prescott (1989), cap.~6, secc.~6.1--6.4}
\nota{Principio de optimalidad, operador de Bellman y teorema de la aplicación
contractiva, a nivel de posgrado. La versión básica de estas herramientas y el
análisis real del cap.~3 de Stokey et al. están en los Bloques 3, 4 y 13 del
sílabo de economía matemática, que se dan por conocidos. El cap.~6 de Stokey et
al., de consulta, prueba la estabilidad de la senda óptima.}

\tema{14.2 Programación dinámica estocástica y lineal-cuadrática}
\lecturas{\AC{16}}
\articulos{Ljungqvist y Sargent (2018), caps.~3 y 5}
\consultalib{Heer y Maußner (2005), secc.~2.1--2.2 \quad Stokey, Lucas y Prescott (1989), cap.~9 (secc.~9.1, 9.2, 9.5 y 9.6) y secc.~10.1 \quad Ljungqvist y Sargent (2018), cap.~2, secc.~2.1--2.4}
\nota{Choques markovianos en la ecuación de Bellman; el regulador
lineal-cuadrático y la ecuación de Riccati. Las cadenas de Markov del cap.~2 de
Ljungqvist y Sargent (y del tema 13.5 de economía matemática) son la base para
discretizar los choques en el tema 14.7.}

\tema{14.3 Equilibrio competitivo, óptimo de Pareto y mercados completos}
\articulos{Ljungqvist y Sargent (2018), cap.~8}
\consultalib{Prácticas UP (2024), secc.~1 \quad Stokey, Lucas y Prescott (1989), caps.~15--16}
\nota{El programa MAT616 remite a los caps.~5 y 13, que en la 4.ª edición son la
programación lineal-cuadrática (ya leída en el tema 14.2) y la valoración de
activos (Bloques 4 y 5 del sílabo de economía financiera). El equilibrio con
mercados completos es el cap.~8, que incluye el problema del planificador con
ponderaciones de Pareto: el método de Negishi.}
\nota{Teoremas del bienestar en economías estáticas y de horizonte infinito;
mercados de Arrow--Debreu frente a mercados secuenciales; economías con
producción. El programa MAT616 lo desarrolla con notas de clase.}
\nota{El capítulo 5 de Acemoglu (tema 4.6) es la versión de este tema para el
crecimiento.}

\tema{14.4 Equilibrio competitivo recursivo}
\articulos{Ljungqvist y Sargent (2018), caps.~7 y 12 \quad Prescott y Mehra (1980)}
\consulta{Mehra (1984)}
\consultalib{Stokey, Lucas y Prescott (1989), caps.~17--18}
\nota{Ley de movimiento agregada y condición de consistencia: el lenguaje de la
macroeconomía cuantitativa moderna.}

\tema{14.5 Soluciones locales: log-linealización y ecuaciones en diferencias}
\articulos{Uhlig (2001)}
\consultalib{\GM{3.5--3.8} \quad Heer y Maußner (2005), secc.~2.3 \quad Sánchez y Maraví (2009)}
\nota{Aproximación alrededor del estado estacionario y cálculo de funciones
impulso-respuesta.}

\tema{14.6 Modelos lineales con expectativas racionales: Blanchard--Kahn, Klein y McCallum}
\articulos{Blanchard y Kahn (1980) \quad Klein (2000)}
\consulta{McCallum (1998) \quad Hamilton (1994), cap.~4}
\nota{Generaliza los métodos del tema 9.1. Usa la diagonalización (tema 2.4 del
sílabo de economía matemática), la factorización de Cholesky (tema 2.1 del de
métodos numéricos) y la descomposición de Schur generalizada (tema 9.1 de ese
mismo sílabo).}
\nota{El apéndice 2.7 de Walsh (tema 8.5) resuelve paso a paso un modelo monetario
con este enfoque.}

\tema{14.7 Métodos numéricos: iteración de la función de valor en Python}
\supl{\KR{10}}
\articulos{Kiusalaas (2013), cap.~1}
\consultalib{Heer y Maußner (2005), caps.~1, 3, 4 y 8 \quad Sánchez y Maraví (2009) \quad Ljungqvist y Sargent (2018), cap.~4}
\nota{Implementación numérica del modelo neoclásico como problema de
programación dinámica (programa 1ECO90). El modelo se usa solo como banco de
pruebas del método: no reabre el Bloque 4.}

\tema{14.8 Calibración y simulación de modelos RBC y DSGE}
\lecturas{\CSV{Apéndices B y C}}
\supl{\KR{10}}
\consultalib{\GM{2.2--2.4}}
\nota{Se estudia dentro del Bloque 12, después del tema 12.7: se calibra y simula
el RBC del tema 12.5 y, al llegar al 12.10, el modelo neokeynesiano. Retoma la
metodología de calibración de Cooley y Prescott (1995).}

%------------------------------------------------------------------------------
\bloque{Bloque 15 --- Fricciones y heterogeneidad en el ciclo económico --- 5 semanas}
{Extensiones del modelo de ciclo que los programas avanzados de la PUCP tratan
como frontera: el mercado laboral con búsqueda, las fricciones financieras y los
mercados incompletos con agentes heterogéneos, con su aplicación a la política
monetaria. Requiere los Bloques 10, 12 y 14; el último tema, también el 13. Es
el destino final del recorrido. Las crisis bancarias (Diamond--Dybvig), las
paradas súbitas y el default soberano se estudian en los sílabos de economía
financiera y de economía internacional.}

\tema{15.1 Búsqueda y emparejamiento en el modelo de ciclo económico real}
\supl{\KR{20}}
\articulos{Merz (1995) \quad Rogerson, Shimer y Wright (2005)}
\consultalib{Ljungqvist y Sargent (2018), secc.~29.2--29.3}
\consulta{Andolfatto (1996) \quad Pissarides (2000)}
\nota{Une el modelo DMP del tema 10.4 con el RBC del 12.5. Eficiencia del
equilibrio y condición de Hosios.}

\tema{15.2 Desempleo y vacantes en el ciclo: el puzzle de Shimer}
\articulos{Shimer (2005)}
\consultalib{Ljungqvist y Sargent (2018), cap.~30}
\nota{El modelo DMP calibrado no genera la volatilidad observada de la razón
vacantes/desempleo.}

\tema{15.3 Informalidad laboral y ciclo económico}
\articulos{Leyva y Urrutia (2020)}
\consulta{Fernández y Meza (2015) \quad Finkelstein Shapiro (2014) \quad Leyva y
Urrutia (2023) \quad Céspedes y Ramírez-Rondán (2021) \quad Jaramillo y Campos
(2021)}
\nota{Regulación laboral, informalidad y fluctuaciones; el mercado laboral en la
pandemia. Incluye evidencia para el Perú.}

\tema{15.4 Fricciones financieras y el acelerador financiero}
\lecturas{\RO{10.1--10.3} \quad \WA{10.5}}
\supl{\KR{19}}
\articulos{Bernanke, Gertler y Gilchrist (1999)}
\consulta{Bernanke y Gertler (1989) \quad Carlstrom y Fuerst (1997)}
\nota{Costos de agencia y patrimonio neto como amplificadores de los choques.}

\tema{15.5 Restricciones de colateral y choques financieros}
\articulos{Kiyotaki y Moore (1997) \quad Quadrini (2011)}
\consulta{Jermann y Quadrini (2012) \quad Kehoe, Midrigan y Pastorino (2018)}
\nota{Ciclos de crédito con colateral; la cuña laboral como síntoma de
fricciones financieras (enlaza con el tema 12.8).}

\tema{15.6 Riesgo idiosincrático, autoseguro y mercados incompletos}
\lecturas{\AC{17.4--17.5}}
\supl{\KR{21}}
\articulos{Aiyagari (1994) \quad Ljungqvist y Sargent (2018), caps.~17 y 18}
\articulosB{Huggett (1993)}
\consultalib{Heer y Maußner (2005), caps.~5--7 \quad Stokey, Lucas y Prescott (1989), caps.~12--13}
\nota{Primer modelo de agentes heterogéneos del plan. Generaliza el ahorro
precautorio del tema 12.4 al equilibrio general: Huggett (1993) lo hace en una
economía de dotaciones con un solo activo, y Aiyagari (1994), con producción y
capital. Los caps.~12 y 13 de Stokey et al., de consulta, garantizan que exista
la distribución estacionaria.}
\nota{Acemoglu presenta el modelo de Bewley y las generaciones solapadas con
incertidumbre.}

\tema{15.7 Agentes heterogéneos con riesgo agregado: Krusell y Smith}
\articulosB{Krusell y Smith (1998)}
\nota{Con riesgo agregado, la distribución de la riqueza entra en el estado del
problema. La agregación aproximada la resume con unos pocos momentos y explica
por qué, en este modelo, los agregados se parecen a los del agente
representativo.}

\tema{15.8 Métodos de solución en el espacio de secuencias}
\articulosB{Auclert, Bardóczy, Rognlie y Straub (2021)}
\nota{Funciones impulso-respuesta de modelos con agentes heterogéneos a partir
de los jacobianos de secuencia: el método estándar para resolver y estimar
modelos HANK, con el paquete \texttt{sequence-jacobian} de Python. Prepara el
tema 15.9.}

\tema{15.9 Política monetaria con agentes heterogéneos}
\articulos{Kaplan, Moll y Violante (2018) \quad Auclert (2019)}
\consulta{Kekre y Lenel (2022) \quad Werning (2015)}
\consultalib{Kaplan y Violante (2018)}
\nota{Canales directos e indirectos de la política monetaria y su efecto
redistributivo. Requiere los temas 13.1, 15.6 y 15.8.}

%------------------------------------------------------------------------------
\bigskip
\section*{Anexo A. De los capítulos a los bloques}
\addcontentsline{toc}{section}{Anexo A. De los capítulos a los bloques}

\subsection*{De Gregorio, \emph{Macroeconomía: Teoría y Políticas}}
{\filasalternas
\noindent\begin{tabularx}{\textwidth}{cXc}
\toprule
\rowcolor{tonoOscuro}\cab{Cap.} & \cab{Título} & \cab{Bloque} \\
\midrule
3  & Consumo & 2 \\
4  & Inversión & 2 \\
5  & El gobierno y la política fiscal & 2 \\
6  & La economía cerrada & 3 \\
7  & Economía abierta: la cuenta corriente & 3 \\
8  & Economía abierta: el tipo de cambio real & 3 \\
9  & Más sobre tipo de cambio real y cuenta corriente & 3 \\
10 & Introducción al crecimiento económico & 4 \\
11 & El modelo neoclásico de crecimiento & 4 \\
12 & Modelos de crecimiento: extensiones & 4--5 \\
13 & Evidencia empírica & 4 \\
14 & Crecimiento económico con ahorro óptimo & 4 \\
15 & Teoría cuantitativa, neutralidad y demanda por dinero & 8 \\
16 & Oferta de dinero, política monetaria e inflación & 8 \\
17 & Política monetaria y mercados financieros & 8 \\
18 & Introducción a las fluctuaciones de corto plazo & 12 \\
19 & El modelo keynesiano de economía cerrada: IS-LM & 1 \\
20 & El modelo de Mundell--Fleming & 1 \\
21 & La oferta agregada y la curva de Phillips & 1 \\
22 & Oferta, demanda agregada y políticas macroeconómicas & 1 \\
23 & Fluctuaciones en modelos del ciclo económico real & 12 \\
24 & Los mercados del trabajo y del crédito en el ciclo & 10 \\
25 & Inconsistencia intertemporal y política monetaria & 9 \\
\bottomrule
\end{tabularx}
}


\bigskip
\subsection*{Heijdra, \emph{Foundations of Modern Macroeconomics} (3.ª ed.)}
{\filasalternas
\noindent\begin{tabularx}{\textwidth}{cXc}
\toprule
\rowcolor{tonoOscuro}\cab{Cap.} & \cab{Título} & \cab{Bloque} \\
\midrule
1  & Review of the AD-AS model & 1.1 \\
2  & The open economy & 1 \\
3  & Dynamics in aggregate demand and supply & 1 \\
4  & Perfect foresight and economic policy & 1--2 \\
5  & Rational expectations and economic policy & 9 \\
6  & The government budget deficit & 2 \\
7  & A closer look at the labour market & 10 \\
8  & Search in the labour market & 10 \\
9  & Dynamic inconsistency in public and private decision making & 9 \\
10 & Money & 8 \\
11 & New Keynesian economics & 11 \\
12 & Exogenous economic growth --- Solow-Swan & 4 \\
13 & Exogenous economic growth --- Ramsey-Cass-Koopmans & 4 \\
14 & Endogenous economic growth & 5 \\
15 & Overlapping generations in continuous time & 7 \\
16 & Overlapping generations in discrete time & 7 \\
17 & Decision making in a stochastic environment & 12 \\
18 & DSGE --- New classical models & 12 \\
19 & DSGE --- New Keynesian models & 12 \\
\bottomrule
\end{tabularx}
}


\bigskip
\subsection*{Romer, \emph{Advanced Macroeconomics} (5.ª ed.)}
\noindent Solo se listan los capítulos utilizados; las secciones excluidas están
en el Anexo D.
\medskip

{\filasalternas
\noindent\begin{tabularx}{\textwidth}{l>{\raggedright\arraybackslash}X>{\raggedright\arraybackslash}p{2.5cm}>{\raggedright\arraybackslash}p{3.2cm}}
\toprule
\rowcolor{tonoOscuro}\cab{Cap.} & \cab{Título} & \cab{Secciones} & \cab{Tema} \\
\midrule
1 & The Solow Growth Model & 1.1, 1.7--1.8 (1.2--1.6, consulta) & 4.1, 4.5, 4.9 \\
2 & Infinite-Horizon and Overlapping-Generations Models & 2.8--2.12 (2.1--2.7, consulta) & 7.1, 7.2, 7.5 \\
3 & Endogenous Growth & 3.4--3.8 (3.1--3.3, consulta) & 5.3, 5.4, 5.7 \\
4 & Cross-Country Income Differences & 4.2--4.6 (4.1, consulta) & 4.5, 6.1, 6.2 \\
5 & Real-Business-Cycle Theory & completo & 8.2, 12.1, 12.5, 12.7 \\
6 & Nominal Rigidity & 6.1--6.3, 6.5--6.9 (6.4, consulta) & 1.4, 9.2, 11.1--11.4 \\
7 & DSGE Models of Fluctuations & completo & 11.5, 12.9--12.11, 13.6 \\
8 & Consumption & 8.1--8.4, 8.6--8.7 & 2.1, 12.3, 12.4 \\
9 & Investment & completo & 2.2, 2.3 \\
10 & Financial Markets and Financial Crises & 10.1--10.3 & 15.4 \\
11 & Unemployment & completo & 10.2, 10.4, 10.5 \\
12 & Monetary Policy & completo & 8.4, 8.8, 9.4, 9.5, 13.1, 13.2, 13.5 \\
13 & Budget Deficits and Fiscal Policy & 13.1--13.8 & 2.4--2.8 \\
\bottomrule
\end{tabularx}
}

\bigskip
\subsection*{Campante, Sturzenegger y Velasco, \emph{Advanced Macroeconomics: An Easy Guide}}
\noindent Solo se listan los capítulos utilizados; los excluidos están en el
Anexo D. En el capítulo 8 queda fuera la subsección sobre la cuenta corriente en
la economía abierta.
\medskip

{\filasalternas
\noindent\begin{tabularx}{\textwidth}{l>{\raggedright\arraybackslash}X>{\raggedright\arraybackslash}p{2.5cm}>{\raggedright\arraybackslash}p{3.2cm}}
\toprule
\rowcolor{tonoOscuro}\cab{Cap.} & \cab{Título} & \cab{Secciones} & \cab{Tema} \\
\midrule
2 & Growth theory preliminaries & consulta & 4.1, 4.2, 4.4 \\
3 & The neoclassical growth model & consulta & 4.7 \\
5 & Endogenous growth models I: Escaping diminishing returns & 5.5 (5.1--5.4, consulta) & 5.1, 6.4 \\
6 & Endogenous growth models II: Technological change & 6.1--6.4 & 5.4, 5.5, 5.7 \\
7 & Proximate and fundamental causes of growth & 7.2 (7.1, consulta) & 4.5, 6.1 \\
8 & Overlapping generations models & 8.1--8.3 & 7.1, 7.2, 7.4 \\
9 & Pension systems and transitions & 9.1--9.2 & 7.3, 7.5 \\
10 & Unified growth theory & completo & 6.5 \\
11 & Consumption & completo & 2.1 \\
12 & Consumption under uncertainty and macro finance & 12.1--12.2 & 12.3, 12.4 \\
13 & Investment & 13.1--13.2 & 2.2, 2.3 \\
14 & Real business cycles & 14.1--14.3 & 12.5--12.7 \\
15 & (New) Keynesian theories of fluctuations & 15.2--15.4 (15.1, consulta) & 1.1, 1.5, 9.2, 11.1, 11.3, 12.10 \\
16 & Unemployment & 16.1--16.5 & 10.1--10.4 \\
17 & Fiscal policy I: Public debt and the effectiveness of fiscal policy & 17.1--17.4, 17.7--17.8 & 2.4--2.6, 11.6 \\
18 & Fiscal policy II: The long-run determinants of fiscal policy & 18.1--18.3 & 2.7, 2.8, 4.8 \\
19 & Monetary policy: An introduction & completo & 8.4--8.6 \\
20 & Rules vs Discretion & completo & 9.4, 9.5 \\
21 & Recent debates in monetary policy & completo & 13.5, 13.6 \\
22 & New developments in monetary and fiscal policy & 22.1--22.3 & 7.6, 8.7, 13.5 \\
Ap.~B, C & Simulating an RBC model; Simulating a DSGE model & completo & 14.8 \\
\bottomrule
\end{tabularx}
}

\bigskip
\subsection*{Acemoglu, \emph{Introduction to Modern Economic Growth}}
\noindent Queda fuera el capítulo 19 (Anexo D). El capítulo 6, las secciones
7.1--7.7 y los apéndices A y B son herramientas matemáticas que se estudian en
el sílabo de economía matemática; los capítulos 22 y 23 son opcionales.
\medskip

{\filasalternas
\noindent\begin{tabularx}{\textwidth}{l>{\raggedright\arraybackslash}X>{\raggedright\arraybackslash}p{2.5cm}>{\raggedright\arraybackslash}p{3.2cm}}
\toprule
\rowcolor{tonoOscuro}\cab{Cap.} & \cab{Título} & \cab{Secciones} & \cab{Tema} \\
\midrule
1 & Economic Growth and Economic Development: The Questions & completo & 4.1 \\
2 & The Solow Growth Model & 2.1--2.7 & 4.2, 4.3 \\
3 & The Solow Model and the Data & 3.1--3.6 & 4.1, 4.4, 4.5 \\
4 & Fundamental Determinants of Differences in Economic Performance & completo & 6.1 \\
5 & Foundations of Neoclassical Growth & completo & 4.6 \\
7 & Review of the Theory of Optimal Control & 7.8 & 2.3 \\
8 & The Neoclassical Growth Model & 8.1--8.10 & 4.7, 4.8 \\
9 & Growth with Overlapping Generations & 9.1--9.8 & 7.1--7.4 \\
10 & Human Capital and Economic Growth & completo & 5.2 \\
11 & First-Generation Models of Endogenous Growth & completo & 5.1 \\
12 & Modeling Technological Change & completo & 5.3 \\
13 & Expanding Variety Models & completo & 5.4 \\
14 & Models of Schumpeterian Growth & completo & 5.5 \\
15 & Directed Technological Change & completo & 5.6 \\
16 & Stochastic Dynamic Programming & completo & 14.2 \\
17 & Stochastic Growth Models & 17.1--17.6 & 6.4, 12.5, 15.6 \\
18 & Diffusion of Technology & completo & 6.2 \\
20 & Structural Change and Economic Growth & completo & 6.3 \\
21 & Structural Transformations and Market Failures in Development & 21.1--21.7 & 6.4, 6.5 \\
22 & Institutions, Political Economy and Growth & opcional & 6.6 \\
23 & Political Institutions and Economic Growth & opcional & 6.7 \\
24 & Epilogue: Mechanics and Causes of Economic Growth & 24.2 (24.1 y 24.3, opcionales) & 6.5, 6.7 \\
Ap.~C & Brief Review of Dynamic Games & opcional & 6.7 \\
\bottomrule
\end{tabularx}
}

\bigskip
\subsection*{Walsh, \emph{Monetary Theory and Policy} (3.ª ed.)}
\noindent Quedan fuera el capítulo 9 y la sección 10.4 (Anexo D).
\medskip

{\filasalternas
\noindent\begin{tabularx}{\textwidth}{l>{\raggedright\arraybackslash}X>{\raggedright\arraybackslash}p{2.5cm}>{\raggedright\arraybackslash}p{3.2cm}}
\toprule
\rowcolor{tonoOscuro}\cab{Cap.} & \cab{Título} & \cab{Secciones} & \cab{Tema} \\
\midrule
1 & Empirical Evidence on Money, Prices, and Output & completo & 8.2 \\
2 & Money-in-the-Utility Function & completo & 8.5, 8.6 \\
3 & Money and Transactions & completo & 8.1, 8.3, 8.5 \\
4 & Money and Public Finance & completo & 8.4, 8.6, 8.7 \\
5 & Money in the Short Run: Informational and Portfolio Rigidities & completo & 9.3 \\
6 & Money in the Short Run: Nominal Price and Wage Rigidities & 6.1--6.3 & 9.3, 11.1--11.3, 11.5, 12.9, 12.10 \\
7 & Discretionary Policy and Time Inconsistency & completo & 9.4, 9.5 \\
8 & New Keynesian Monetary Economics & completo & 12.10, 13.1, 13.3 \\
10 & Financial Markets and Monetary Policy & 10.1--10.3, 10.5--10.8 & 8.8, 13.2, 13.5, 15.4 \\
11 & Monetary Policy Operating Procedures & completo & 8.9 \\
\bottomrule
\end{tabularx}
}

\bigskip
\subsection*{Galí, \emph{Monetary Policy, Inflation, and the Business Cycle} (1.ª ed.)}
\noindent Queda fuera el capítulo 7 (Anexo D).
\medskip

{\filasalternas
\noindent\begin{tabularx}{\textwidth}{l>{\raggedright\arraybackslash}X>{\raggedright\arraybackslash}p{2.5cm}>{\raggedright\arraybackslash}p{3.2cm}}
\toprule
\rowcolor{tonoOscuro}\cab{Cap.} & \cab{Título} & \cab{Secciones} & \cab{Tema} \\
\midrule
1 & Introduction & completo & 12.10 \\
2 & A Classical Monetary Model & completo & 8.5 \\
3 & The Basic New Keynesian Model & completo & 12.10 \\
4 & Monetary Policy Design in the Basic New Keynesian Model & completo & 13.1 \\
5 & Monetary Policy Tradeoffs: Discretion versus Commitment & completo & 13.3 \\
6 & A Model with Sticky Wages and Prices & completo & 13.4 \\
8 & Main Lessons and Some Extensions & completo & 13.6 \\
\bottomrule
\end{tabularx}
}

\bigskip
\subsection*{Galindo y Montecinos, \emph{Macrodinámica, Tomo I} (consulta)}
\noindent Libro de consulta: no entra en el cómputo de horas. Solo se listan las
partes seleccionadas; quedan fuera el capítulo 4 (versión intermedia del modelo
con trabajo constante) y las secciones de códigos.
\medskip

{\filasalternas
\noindent\begin{tabularx}{\textwidth}{l>{\raggedright\arraybackslash}X>{\raggedright\arraybackslash}p{2.5cm}>{\raggedright\arraybackslash}p{3.2cm}}
\toprule
\rowcolor{tonoOscuro}\cab{Cap.} & \cab{Título} & \cab{Secciones} & \cab{Tema} \\
\midrule
1 & Panorama de los modelos RBC & 1.2--1.4 & 12.1, 12.5 \\
2 & Fundamentos de Dynare & 2.2--2.4 & 14.8 \\
3 & Modelo RBC con solución analítica & 3.5--3.8 & 14.5 \\
5 & Modelo RBC con trabajo variable & completo & 12.5 \\
6 & Modelo RBC con choque a la inversión y utilización variable del capital & completo & 12.6 \\
\bottomrule
\end{tabularx}
}

\bigskip
\subsection*{Azzimonti, Krusell, McKay y Mukoyama, \emph{Macroeconomics} (suplemento)}
\noindent Solo se listan los capítulos que entran en las horas. Los capítulos 3,
4, 5, 13 y 22 quedan como consulta en los Bloques 4 y 5, y el capítulo 9
(técnicas en tiempo continuo) corresponde al sílabo de economía matemática.
\medskip

{\filasalternas
\noindent\begin{tabularx}{\textwidth}{cXc}
\toprule
\rowcolor{tonoOscuro}\cab{Cap.} & \cab{Título} & \cab{Tema} \\
\midrule
6 & Welfare & 7.2 \\
7 & Uncertainty & 12.2, 12.3 \\
8 & Empirical Strategies and Quantitative Macroeconomics & 12.7 \\
10 & Computational Tools & 14.7, 14.8 \\
11 & Consumption & 2.1, 12.3 \\
12 & Labor Supply & 12.5 \\
14 & Real Business Cycles & 12.5, 12.7 \\
15 & Government and Public Policies & 2.4, 2.6 \\
17 & Money & 8.3, 8.5 \\
18 & Nominal Frictions and Business Cycles & 11.2, 12.10 \\
19 & Credit Market Frictions & 10.1, 15.4 \\
20 & Frictional Labor Markets & 10.4, 15.1 \\
21 & Inequality & 15.6 \\
23 & International Macroeconomics & 3.2, 3.4 \\
25 & Sustainability & 4.9 \\
\bottomrule
\end{tabularx}
}

%------------------------------------------------------------------------------
\newpage
\section*{Anexo B. Temas cubiertos por un solo texto}
\addcontentsline{toc}{section}{Anexo B. Temas cubiertos por un solo texto}

Conviene tener presente dónde \emph{no} hay redundancia entre De Gregorio y
Heijdra, para no saltarse material que carece de sustituto.

\medskip
\noindent\destacado{Solo en De Gregorio:}
\begin{itemize}[leftmargin=1.4em,itemsep=2pt]
  \item Cuenta corriente intertemporal y tipo de cambio real (\DG{7}, \DG{8}, \DG{9}).
  \item Evidencia empírica del crecimiento (\DG{13}).
  \item Política monetaria y mercados financieros (\DG{17}).
  \item Evidencia y contexto de economías emergentes y de América Latina
        (transversal a todo el libro).
\end{itemize}

\medskip
\noindent\destacado{Solo en Heijdra:}
\begin{itemize}[leftmargin=1.4em,itemsep=2pt]
  \item Previsión perfecta y análisis de sendero de silla (\HJ{4}).
  \item Métodos de solución con expectativas racionales (\HJ{5}).
  \item Búsqueda y emparejamiento en el mercado laboral (\HJ{8}).
  \item Fundamentos microeconómicos de la nueva economía keynesiana (\HJ{11}).
  \item Generaciones solapadas en tiempo continuo y discreto (\HJ{15}, \HJ{16}).
  \item Programación dinámica y decisiones estocásticas (\HJ{17}).
  \item Modelos DSGE resueltos, clásicos y neokeynesianos (\HJ{18}, \HJ{19}).
\end{itemize}


%------------------------------------------------------------------------------
\bigskip
\section*{Anexo C. Capítulos del suplemento excluidos del plan}
\addcontentsline{toc}{section}{Anexo C. Capítulos del suplemento excluidos}

Los siguientes capítulos de \textcolor{colorK}{\textbf{K}} \emph{no} se han
incorporado, porque abren materia que este sílabo no contempla. No son un olvido:
quedan disponibles como continuación natural una vez cerrado el Bloque 15.

\medskip
{\filasalternas
\noindent\begin{tabularx}{\textwidth}{cXX}
\toprule
\rowcolor{tonoOscuro}\cab{Cap.} & \cab{Título} & \cab{Motivo de la exclusión} \\
\midrule
16 & Asset Prices & Valoración de activos; se estudia en el sílabo de economía financiera. \\
24 & Sovereign Debt and Default Risk & Adyacente al tema 2.6, pero es una literatura distinta de la sostenibilidad de la deuda. \\
\bottomrule
\end{tabularx}
}

\medskip
\noindent\destacado{Una salvedad sobre el Bloque 9.} Expectativas racionales,
crítica de Lucas e inconsistencia dinámica no reciben suplemento: ningún capítulo
de \textcolor{colorK}{\textbf{K}} está dedicado a ese material. \HJ{5} y \HJ{9}
siguen siendo la referencia principal; Walsh, Romer y Campante et al. los
complementan.


%------------------------------------------------------------------------------
\newpage
\section*{Anexo D. Capítulos y secciones que no entran en el plan}
\addcontentsline{toc}{section}{Anexo D. Capítulos y secciones excluidos}

Quedan fuera los contenidos de macroeconomía internacional y de finanzas y
banca, que tienen sílabos propios, y el material introductorio. La
asignación se hizo a partir de los índices de cada libro.

\medskip
{\filasalternas
\noindent\begin{tabularx}{\textwidth}{lXX}
\toprule
\rowcolor{tonoOscuro}\cab{Cap. / secc.} & \cab{Contenido} & \cab{Motivo de la exclusión} \\
\midrule
\rowcolor{tonoHielo}\multicolumn{3}{l}{\textcolor{colorR}{\textbf{Romer}}} \\
Intro. & Introduction & Material introductorio. \\
8.5 & Consumption and Risky Assets & Valoración de activos; coherente con la exclusión de \KR{16}. \\
10.4--10.8 & Volatilidad excesiva, modelo de Diamond--Dybvig, contagio y crisis financieras & Economía financiera y bancaria. \\
13.9 & A Model of Sovereign Debt Crises & Macroeconomía internacional; coherente con la exclusión de \KR{24}. \\
\midrule
\rowcolor{tonoHielo}\multicolumn{3}{l}{\textcolor{colorCSV}{\textbf{Campante, Sturzenegger y Velasco}}} \\
1 & Introduction & Material introductorio. \\
4 & An application: The small open economy & Macroeconomía internacional. \\
8.3 (subsecc.) & Revisiting the current account in the open economy & Macroeconomía internacional. \\
12.3 & Macroeconomics and finance (CAPM de consumo, prima por riesgo) & Economía financiera; coherente con la exclusión de \KR{16}. \\
13.3 & Investment in the open economy & Macroeconomía internacional. \\
\midrule
\rowcolor{tonoHielo}\multicolumn{3}{l}{\textcolor{colorW}{\textbf{Walsh}}} \\
10.4 & Macrofinance & Valoración de activos; coherente con la exclusión de \KR{16}. \\
\bottomrule
\end{tabularx}
}

\medskip
\noindent Los capítulos de economía abierta de Acemoglu (\AC{19}), Walsh
(\WA{9}) y Galí (\GA{7}) entran como consulta, sin horas, en los temas 6.2 y
13.1. Su estudio con horas corresponde al sílabo de economía internacional
(\GA{7}, en el tema 8.4 de ese plan).

\medskip
\noindent Las herramientas matemáticas tampoco entran: el apéndice matemático de
Heijdra, el apéndice A de Campante et al., el capítulo 6, las secciones 7.1--7.7
y los apéndices A y B de Acemoglu y el capítulo 9 del suplemento se estudian en
el sílabo de economía matemática.

%------------------------------------------------------------------------------
\newpage
\section*{Anexo E. Programas de la PUCP utilizados}
\addcontentsline{toc}{section}{Anexo E. Programas de la PUCP utilizados}

Se usó la versión más reciente de cada programa. Los números entre paréntesis
remiten a los temas de este sílabo.

\medskip
{\small\filasalternas
\noindent\begin{tabularx}{\textwidth}{>{\raggedright\arraybackslash}p{3.1cm}>{\raggedright\arraybackslash}X>{\raggedright\arraybackslash}X}
\toprule
\rowcolor{tonoOscuro}\cab{Programa} & \cab{Incorporado} & \cab{No incorporado} \\
\midrule
\textbf{1ECO74} (2026-1)\newline Macroeconomía Intermedia\newline\emph{Maestría en Economía}
 & Consumo e inversión; crecimiento exógeno y endógeno; RBC en economía cerrada y sus extensiones; fundamentos de la teoría monetaria, crítica de Lucas, señoreaje e inflación; modelo neokeynesiano y diseño de la política monetaria (2.1, 2.3, 4.1--4.3, 4.7--4.9, 5.1, 6.1, 6.2, 6.4, 6.5, 7.1, 8.4, 8.5, 9.2, 12.1, 12.4--12.7, 12.10, 12.11, 13.1--13.3, 15.4).
 & Revisión histórica de la economía peruana; RBC en economía abierta y términos de intercambio; dolarización y políticas macroprudenciales; Campbell (1999) y Mehra y Prescott (1985), de valoración de activos; Freixas y Rochet (2008), de banca. \\[3pt]
\textbf{2ECO03} (2025-2)\newline Macroeconomía Avanzada\newline\emph{Maestría en Economía}
 & Búsqueda y emparejamiento en el mercado laboral (módulo 4) y artículos clásicos del acelerador financiero en economía cerrada (módulo 1) (10.4, 15.1, 15.2, 15.4, 15.5).
 & Módulo 1 (fricciones financieras y política macroprudencial en economías abiertas), módulo 2 (modelo neokeynesiano de economía abierta) y módulo 3 (dolarización financiera). Su sumilla describe, por error, un curso de microeconomía. \\[3pt]
\textbf{1ECO90} (2026-1)\newline Temas de Macroeconomía Avanzada\newline\emph{Maestría en Economía}
 & Métodos numéricos; mercado laboral y ciclo económico; formación de precios; determinantes de la inflación y su control (8.4, 8.7, 10.4, 11.5, 12.2, 12.5, 12.6, 13.2, 14.1, 14.7, 15.1--15.3).
 & La versión 2025-1 no se usó. \\[3pt]
\textbf{1ECO14} (2026-1, horario 722)\newline Economía Monetaria
 & Naturaleza y demanda del dinero, dinero en equilibrio general, dinero y déficit público, rigideces nominales, mercados financieros y transmisión, inconsistencia dinámica, metas de inflación y procedimientos operativos (8.1--8.6, 8.8, 8.9, 9.3--9.5, 11.1, 12.9, 12.10, 13.1, 13.2).
 & Transmisión en economía abierta, política macroprudencial y política monetaria del BCRP, con sus lecturas; Freixas y Rochet (2023), de banca; Christiano e Ikeda (2016) y Jeanne y Korinek (2010). Se usó el horario 722; el 721, del mismo semestre, es casi idéntico. \\[3pt]
\textbf{ECO324} (2026-1)\newline Tópicos de Macroeconomía Avanzada\newline\emph{Pregrado}
 & Modelo RBC canónico, calibración, log-linealización, funciones impulso-respuesta y contabilidad del ciclo; modelos neokeynesianos; lecturas sobre trampa de liquidez, \emph{forward guidance} y HANK (11.1, 11.6, 12.1, 12.5, 12.7, 12.8, 12.10, 13.1, 13.5, 13.6, 14.5, 15.9).
 & Farhi y Werning (2016), de política macroprudencial. La versión 2025-1 no se usó. \\[3pt]
\textbf{ECO338} (2026-2)\newline Tópicos de Macroeconomía Avanzada~2\newline\emph{Pregrado}
 & RBC y extensiones, modelo neokeynesiano básico y fricciones de crédito (secciones 1 a 3) (9.2, 12.1, 12.2, 12.5--12.8, 12.10, 12.11, 13.1, 13.2, 15.4--15.6).
 & Sección 4 (macroeconomía internacional). \\[3pt]
\textbf{MAT616} (2023-2)\newline Macroeconomía Dinámica\newline\emph{Maestría en Matemáticas Aplicadas}
 & Herramientas numéricas, modelo de Solow, programación dinámica, soluciones locales, bienestar y mercados completos, equilibrio competitivo recursivo, crecimiento endógeno de primera generación, calibración y mercados incompletos (4.1, 4.2, 5.1, 12.1, 14.1--14.8, 15.6).
 & Nada; el problema del consumidor (tema 3 del programa) ya está cubierto por el tema 2.1. \\
\bottomrule
\end{tabularx}
}

\medskip
\noindent\destacado{Programas no utilizados}, por corresponder a áreas que siguen
fuera del plan: Finanzas Internacionales (1ECO93 y 1FIN27); Políticas
Macroeconómicas en el Perú (1ECO46); Economía de las Instituciones Financieras
(1ECO86); y Economía Financiera y Bancaria (1ECO87 y 1FIN30).

\medskip
\noindent\destacado{Sobre las ediciones de Galí.} 1ECO14 y 1ECO74 citan la
1.ª edición (2008), que es la que usa el plan; ECO324, ECO338 y 2ECO03 citan la
2.ª (2015).

%------------------------------------------------------------------------------
\bigskip
\section*{Anexo F. Referencias tomadas de los programas de la PUCP}
\addcontentsline{toc}{section}{Anexo F. Referencias de los programas de la PUCP}

Datos bibliográficos según los programas; entre corchetes, el tema donde se
usan. Las referencias de capítulo de Ljungqvist y Sargent, Stokey et al.,
Woodford y Benigno son las que indican los programas, salvo en el tema 14.3,
donde se corrigió la referencia de Ljungqvist y Sargent a la 4.ª edición.

\medskip
{\small
\begin{itemize}[leftmargin=1.4em,itemsep=1pt,label={}]
\item Aiyagari, S. R. (1994). Uninsured Idiosyncratic Risk and Aggregate Saving. \emph{Quarterly Journal of Economics}, 109(3), 659--684. [15.6]
\item Andolfatto, D. (1996). Business Cycles and Labor-Market Search. \emph{American Economic Review}. [15.1]
\item Auclert, A. (2019). Monetary Policy and the Redistribution Channel. \emph{American Economic Review}, 109(6), 2333--2367. [15.9]
\item Balboni, C., O. Bandiera, R. Burgess, M. Ghatak y A. Heil (2021). Why Do People Stay Poor? \emph{Quarterly Journal of Economics}. [6.4]
\item Barro, R. J. y D. B. Gordon (1983). A Positive Theory of Monetary Policy in a Natural Rate Model. \emph{Journal of Political Economy}, 91(4), 589--610. [9.4]
\item Barro, R. J. y X. Sala-i-Martin (2003). \emph{Economic Growth}, 2.ª ed. MIT Press. [4.2]
\item Baumol, W. (1952). The Transactions Demand for Cash: An Inventory Theoretic Approach. \emph{Quarterly Journal of Economics}, 66, 545--556. [8.3]
\item Benhabib, J., R. Rogerson y R. Wright (1991). Homework in Macroeconomics: Household Production and Aggregate Fluctuations. \emph{Journal of Political Economy}. [12.6]
\item Benigno, P. (2025). \emph{Monetary Economics and Policy: A Foundation for Modern Currency Systems}. Princeton University Press. [8.1, 8.5, 8.8]
\item Bernanke, B. y F. Mishkin (1997). Inflation Targeting: A New Framework for Monetary Policy? \emph{Journal of Economic Perspectives}, 11(2), 97--116. [9.5]
\item Bernanke, B. y M. Gertler (1989). Agency Costs, Net Worth, and Business Fluctuations. \emph{American Economic Review}, 79(1), 14--31. [15.4]
\item Bernanke, B. y M. Gertler (1995). Inside the Black Box: The Credit Channel of Monetary Policy Transmission. \emph{Journal of Economic Perspectives}, 9(4), 27--48. [8.8]
\item Bernanke, B., M. Gertler y S. Gilchrist (1999). The Financial Accelerator in a Quantitative Business Cycle Framework. En J. B. Taylor y M. Woodford (eds.), \emph{Handbook of Macroeconomics}, vol.~1, cap.~21, 1341--1393. [15.4]
\item Bindseil, U. (2014). \emph{Monetary Policy Operations and the Financial System}, caps.~1--3. [8.9]
\item BIS (2018). \emph{Central Bank Digital Currencies}. Committee on Payments and Market Infrastructures. [8.1]
\item BIS (2025). The Next-Generation Monetary and Financial System. \emph{BIS Annual Economic Report 2025}, cap.~III, 77--113. [8.1]
\item Blanchard, O. J. y C. M. Kahn (1980). The Solution of Linear Difference Models under Rational Expectations. \emph{Econometrica}, 48(5), 1305--1311. [14.6]
\item Blanchard, O. J. y N. Kiyotaki (1987). Monopolistic Competition and the Effects of Aggregate Demand. \emph{American Economic Review}, 647--666. [11.1]
\item Bloom, D. E., M. Kuhn y K. Prettner (2022). Modern Infectious Diseases: Macroeconomic Impacts and Policy Responses. \emph{Journal of Economic Literature}, 60(1), 85--131. [6.1]
\item Böhme, R., N. Christin, B. Edelman y T. Moore (2015). Bitcoin: Economics, Technology, and Governance. \emph{Journal of Economic Perspectives}, 29(2), 213--238. [8.1]
\item Bueno, G. (2024). Inflation and Price Dynamics Analysis using High-Frequency Data from Supermarkets: Evidence from Perú. Manuscrito. [11.5]
\item Calvo, G. A. (1983). Staggered Prices in a Utility-Maximizing Framework. \emph{Journal of Monetary Economics}, 12, 383--398. [12.10]
\item Carlstrom, C. y T. Fuerst (1997). Agency Costs, Net Worth, and Business Fluctuations: A Computable General Equilibrium Analysis. \emph{American Economic Review}, 87(5), 893--910. [15.4]
\item Carroll, C. (1997). Buffer-Stock Saving and the Life Cycle/Permanent Income Hypothesis. \emph{Quarterly Journal of Economics}, 112(1), 1--55. [12.4]
\item Castillo-Martinez, L. y R. Reis (2026). How Do Central Banks Control Inflation? A Guide for the Perplexed. \emph{Journal of Economic Literature}, 64(1), 195--245. [13.2]
\item Cavallo, A. (2018). More Amazon Effects: Online Competition and Pricing Behaviors. NBER Working Paper 25138. [11.5]
\item Cavallo, A., F. Lippi y K. Miyahara (2024). Large Shocks Travel Fast. \emph{American Economic Review: Insights}, 6(4), 558--574. [11.5]
\item Céspedes, N. y N. Ramírez-Rondán (2021). Job Finding and Separation Rates in an Economy with High Labor Informality. \emph{Research in Labor Economics}. [15.3]
\item Chari, V. V., P. J. Kehoe y E. R. McGrattan (2007). Business Cycle Accounting. \emph{Econometrica}, 75(3), 781--836. [12.8]
\item Clarida, R., J. Galí y M. Gertler (1999). The Science of Monetary Policy: A New Keynesian Perspective. \emph{Journal of Economic Literature}, 37, 1661--1707. [13.1]
\item Cooley, T. F. y E. C. Prescott (1995). Economic Growth and Business Cycles. En T. F. Cooley (ed.), \emph{Frontiers of Business Cycle Research}, Princeton University Press, cap.~1, 1--38. [12.1]
\item Correia, I., E. Farhi, J. P. Nicolini y P. Teles (2013). Unconventional Fiscal Policy at the Zero Bound. \emph{American Economic Review}, 103(4), 1172--1211. [13.5]
\item Davis, S. y J. Haltiwanger (1992). Gross Job Creation, Gross Job Destruction, and Employment Reallocation. \emph{Quarterly Journal of Economics}. [10.4]
\item Dornbusch, R. (1992). Lessons from Experience with High Inflation. \emph{World Bank Economic Review}, 13--32. [8.4]
\item Duncan, R. (2022). Lecture Notes on Economic Growth. Ohio University. [4.1]
\item Duncan, R. y R. Fuentes (2006). Regional Convergence in Chile: New Tests, Old Results. \emph{Cuadernos de Economía}, 43(127), 81--112. [4.3]
\item Eggertsson, G. B. y M. Woodford (2004). Optimal Monetary and Fiscal Policy in a Liquidity Trap. \emph{NBER International Seminar on Macroeconomics 2004}, 75--144. [13.5]
\item Fernández, A. y F. Meza (2015). Informal Employment and Business Cycles in Emerging Economies: The Case of Mexico. \emph{Review of Economic Dynamics}. [15.3]
\item Fernández-Villaverde, J. y D. Sanches (2024). Price-Level Determination under the Gold Standard. Documento de trabajo, SSRN 4744246. [13.2]
\item Finkelstein Shapiro, A. (2014). Self-Employment and Business Cycle Persistence: Does the Composition of Employment Matter for Economic Recoveries? \emph{Journal of Economic Dynamics and Control}. [15.3]
\item Fischer, S. (1977). Long-Term Contracts, Rational Expectations, and the Optimal Money Supply Rule. \emph{Journal of Political Economy}, 85, 191--205. [12.9]
\item Galí, J. (1999). Technology, Employment, and the Business Cycle: Do Technology Shocks Explain Aggregate Fluctuations? \emph{American Economic Review}, 249--271. [12.7]
\item Galí, J., J. D. López-Salido y J. Vallés (2007). Understanding the Effects of Government Spending on Consumption. \emph{Journal of the European Economic Association}, 5(1), 227--270. [11.6]
\item Goodfriend, M. y R. G. King (1997). The New Neoclassical Synthesis and the Role of Monetary Policy. \emph{NBER Macroeconomics Annual}, 12, 231--283. [12.10]
\item Goodhart, C. (1994). What Should Central Banks Do? What Should Be Their Macroeconomic Objectives and Operations? \emph{Economic Journal}. [8.9]
\item Gorodnichenko, Y. y O. Talavera (2017). Price Setting in Online Markets: Basic Facts, International Comparisons, and Cross-Border Integration. \emph{American Economic Review}, 107(1), 249--282. [11.5]
\item Greenwood, J., Z. Hercowitz y G. W. Huffman (1988). Investment, Capacity Utilization, and the Real Business Cycle. \emph{American Economic Review}, 78, 402--417. [12.6]
\item Hamilton, J. D. (1994). \emph{Time Series Analysis}. Princeton University Press, cap.~4. [14.6]
\item Hansen, G. D. (1985). Indivisible Labor and the Business Cycle. \emph{Journal of Monetary Economics}, 16, 309--327. [12.6]
\item Hayashi, F. (1982). Tobin's Marginal q and Average q: A Neoclassical Interpretation. \emph{Econometrica}, 50, 213--224. [2.3]
\item Jaramillo, M. y D. Campos (2021). \emph{La dinámica del mercado laboral peruano: creación y destrucción de empleos y flujos de trabajadores}. GRADE. [15.3]
\item Jermann, U. y V. Quadrini (2012). Macroeconomic Effects of Financial Shocks. \emph{American Economic Review}, 102, 238--271. [15.5]
\item Jones, C. I. (2016). The Facts of Economic Growth. En \emph{Handbook of Macroeconomics}, vol.~2, 3--69. [4.1]
\item Kaplan, G., B. Moll y G. L. Violante (2018). Monetary Policy According to HANK. \emph{American Economic Review}, 108(3), 697--743. [15.9]
\item Kehoe, P. J., V. Midrigan y E. Pastorino (2018). Evolution of Modern Business Cycle Models: Accounting for the Great Recession. \emph{Journal of Economic Perspectives}, 32, 141--166. [15.5]
\item Kekre, R. y M. Lenel (2022). Monetary Policy, Redistribution, and Risk Premia. \emph{Econometrica}, 90(5), 2249--2282. [15.9]
\item King, M. (1999). Challenges for Monetary Policy: New and Old. En \emph{New Challenges for Monetary Policy}, simposio de 1999. [9.5]
\item King, R. G. y S. T. Rebelo (1999). Resuscitating Real Business Cycles. En J. B. Taylor y M. Woodford (eds.), \emph{Handbook of Macroeconomics}, vol.~1, 927--1007. [12.5]
\item King, R. G., C. I. Plosser y S. T. Rebelo (1988). Production, Growth and Business Cycles: I. The Basic Neoclassical Model. \emph{Journal of Monetary Economics}, 21(2--3), 195--232. [12.5]
\item Kiusalaas, J. (2013). \emph{Numerical Methods in Engineering with Python 3}, 3.ª ed. Cambridge University Press, cap.~1. [14.7]
\item Kiyotaki, N. y J. Moore (1997). Credit Cycles. \emph{Journal of Political Economy}, 105(2), 211--248. [15.5]
\item Kiyotaki, N. y J. Moore (2002). Evil Is the Root of All Money. \emph{American Economic Review}, 92(2), 62--66. [8.1]
\item Kiyotaki, N. y R. Wright (1989). On Money as a Medium of Exchange. \emph{Journal of Political Economy}, 97(4). [8.1]
\item Klein, P. (2000). Using the Generalized Schur Form to Solve a Multivariate Linear Rational Expectations Model. \emph{Journal of Economic Dynamics and Control}, 24(10), 1405--1423. [14.6]
\item Klenow, P. J. y B. A. Malin (2010). Microeconomic Evidence on Price-Setting. En \emph{Handbook of Monetary Economics}, vol.~3, 231--284. [11.5]
\item Kydland, F. E. y E. C. Prescott (1982). Time to Build and Aggregate Fluctuations. \emph{Econometrica}, 50, 1345--1370. [12.5]
\item León, D. (1999). La información contenida en los agregados monetarios. \emph{Revista de Estudios Económicos}, 3, Banco Central de Reserva del Perú. [8.3]
\item Lewis, M. y P. Mizen (2000). \emph{Monetary Economics}. Oxford University Press. [8.1, 8.4, 8.9, 9.5]
\item Leyva, G. y C. Urrutia (2020). Informality, Labor Regulation, and the Business Cycle. \emph{Journal of International Economics}. [15.3]
\item Leyva, G. y C. Urrutia (2023). Informal Labor Markets in Times of Pandemic. \emph{Review of Economic Dynamics}. [15.3]
\item Ljungqvist, L. y T. J. Sargent (2018). \emph{Recursive Macroeconomic Theory}, 4.ª ed. MIT Press. [14.2--14.4, 15.6]
\item Long, J. B. y C. I. Plosser (1983). Real Business Cycles. \emph{Journal of Political Economy}, 91, 39--69. [12.5]
\item Lucas, R. E. (1976). Econometric Policy Evaluation: A Critique. \emph{Carnegie-Rochester Conference Series on Public Policy}, 1, 19--46. [9.2]
\item Lucas, R. E. (1994). The Welfare Costs of Inflation. CEPR Publication No.~394, Stanford University. [8.6]
\item Lucas, R. E. (1996). Nobel Lecture: Monetary Neutrality. \emph{Journal of Political Economy}, 104(4), 661--682. [8.6]
\item Lucas, R. E. (2013). \emph{Collected Papers on Monetary Theory}. Harvard University Press. [8.5]
\item Mankiw, N. G. y M. Weinzierl (2006). Dynamic Scoring: A Back-of-the-Envelope Guide. \emph{Journal of Public Economics}, 90(8--9), 1415--1433. [4.8]
\item Mankiw, N. G. y R. Reis (2010). Imperfect Information and Aggregate Supply. En \emph{Handbook of Monetary Economics}, vol.~3, 183--229. [9.3]
\item McCallum, B. T. (1998). Solutions to Linear Rational Expectations Models: A Compact Exposition. \emph{Economics Letters}, 61(2), 143--147. [14.6]
\item McCallum, B. T. (2015). The Bitcoin Revolution. \emph{Cato Journal}, 35(2). [8.1]
\item McCallum, B. T. y E. Nelson (2010). Money and Inflation: Some Critical Issues. En \emph{Handbook of Monetary Economics}, vol.~3, 97--153. [8.5]
\item McKay, A., E. Nakamura y J. Steinsson (2016). The Power of Forward Guidance Revisited. \emph{American Economic Review}, 106(10), 3133--3158. [13.6]
\item Mehra, R. (1984). Recursive Competitive Equilibrium: A Parametric Example. \emph{Economics Letters}, 16(3--4), 273--278. [14.4]
\item Merz, M. (1995). Search in the Labor Market and the Real Business Cycle. \emph{Journal of Monetary Economics}. [15.1]
\item Midrigan, V. (2011). Menu Costs, Multiproduct Firms, and Aggregate Fluctuations. \emph{Econometrica}, 79(4), 1139--1180. [11.5]
\item Mishkin, F. (1995). Symposium on the Monetary Transmission Mechanism. \emph{Journal of Economic Perspectives}, 9(4), 3--10. [8.8]
\item Mishkin, F. (2000). Inflation Targeting in Emerging-Market Countries. \emph{American Economic Review}, 90(2), 105--109. [9.5]
\item Mishkin, F. (2021). \emph{The Economics of Money, Banking and Financial Markets}, 13.ª ed. [8.1, 8.3, 8.8, 8.9, 9.5]
\item Mitchell, W. C. (1913). \emph{Business Cycles}, vol.~3. University of California Press. [12.1]
\item Mortensen, D. y C. Pissarides (1994). Job Creation and Job Destruction in the Theory of Unemployment. \emph{Review of Economic Studies}. [10.4]
\item Nakamura, E. y J. Steinsson (2008). Five Facts about Prices: A Reevaluation of Menu Cost Models. \emph{Quarterly Journal of Economics}, 123(4), 1415--1464. [11.5]
\item Nakamura, E. y J. Steinsson (2013). Price Rigidity: Microeconomic Evidence and Macroeconomic Implications. \emph{Annual Review of Economics}, 5(1), 133--163. [11.5]
\item NobelPrize.org (2018). Scientific Background: Economic Growth, Technological Change, and Climate Change. [4.9]
\item Ogilvie, S. y A. W. Carus (2014). Institutions and Economic Growth in Historical Perspective. En \emph{Handbook of Economic Growth}, vol.~2A, cap.~8. [6.1]
\item Pissarides, C. (2000). \emph{Equilibrium Unemployment Theory}, 2.ª ed. MIT Press. [10.4, 15.1]
\item Poole, W. (1970). Optimal Choice of Monetary Policy Instruments in a Simple Stochastic Macro Model. \emph{Quarterly Journal of Economics}, 84(2), 197--216. [8.9]
\item Prescott, E. C. (1986). Theory Ahead of Business Cycle Measurement. \emph{Federal Reserve Bank of Minneapolis Quarterly Review}, 10, 9--22. [12.7]
\item Prescott, E. C. (2016). RBC Methodology and the Development of Aggregate Economic Theory. NBER Working Paper 22422. [12.7]
\item Prescott, E. C. y R. Mehra (1980). Recursive Competitive Equilibrium: The Case of Homogeneous Households. \emph{Econometrica}, 48(6), 1365--1379. [14.4]
\item Quadrini, V. (2011). Financial Frictions in Macroeconomic Fluctuations. \emph{Economic Quarterly}, Federal Reserve Bank of Richmond, 97, 209--254. [15.5]
\item Quispe, Z. (1998). Una aproximación a la demanda de los principales agregados monetarios en el Perú: junio 1991--mayo 1997. \emph{Revista de Estudios Económicos}, 3, Banco Central de Reserva del Perú. [8.3]
\item Reis, R. (2013). Central Bank Design. \emph{Journal of Economic Perspectives}, 27(4), 17--44. [9.5]
\item Rogerson, R. (1988). Indivisible Labor, Lotteries and Equilibrium. \emph{Journal of Monetary Economics}, 21, 3--16. [12.6]
\item Rogerson, R., R. Shimer y R. Wright (2005). Search-Theoretic Models of the Labor Market: A Survey. \emph{Journal of Economic Literature}, 43, 959--988. [15.1]
\item Rogoff, K. (1985). The Optimal Degree of Commitment to an Intermediate Monetary Target. \emph{Quarterly Journal of Economics}, 100(4). [9.5]
\item Romer, P. (1996). Why, Indeed, in America? Theory, History, and the Origins of Modern Economic Growth. \emph{American Economic Review}, 86(2), 202--206. [6.5]
\item Shimer, R. (2005). The Cyclical Behavior of Equilibrium Unemployment and Vacancies. \emph{American Economic Review}, 95(1), 25--49. [15.2]
\item Sims, C. A. (2010). Rational Inattention and Monetary Economics. En \emph{Handbook of Monetary Economics}, vol.~3, 155--181. [9.3]
\item Smets, F. y R. Wouters (2007). Shocks and Frictions in US Business Cycles: A Bayesian DSGE Approach. \emph{American Economic Review}, 97, 586--606. [12.11]
\item Solow, R. M. (1956). A Contribution to the Theory of Economic Growth. \emph{Quarterly Journal of Economics}, 70(1), 65--94. [4.2]
\item Stokey, N. L. y R. E. Lucas, con E. C. Prescott (1989). \emph{Recursive Methods in Economic Dynamics}. Harvard University Press, caps.~1, 2, 4 y 5. [14.1]
\item Syverson, C. (2011). What Determines Productivity? \emph{Journal of Economic Literature}, 49(2), 326--365. [6.2]
\item Taylor, J. B. (1993). Discretion versus Policy Rules in Practice. \emph{Carnegie-Rochester Conference Series on Public Policy}, 39, 195--214. [13.2]
\item Tobin, J. (1956). The Interest-Elasticity of Transactions Demand for Cash. \emph{Review of Economics and Statistics}, 38, 241--247. [8.3]
\item Tobin, J. (1958). Liquidity Preference as Behavior Towards Risk. \emph{Review of Economic Studies}. [8.3]
\item Uhlig, H. (2001). A Toolkit for Analysing Nonlinear Dynamic Stochastic Models Easily. En R. Marimon y A. Scott (eds.), \emph{Computational Methods for the Study of Dynamic Economies}, Oxford University Press, cap.~3, 30--61. [14.5]
\item Urrutia, C. (1996). Notas sobre crecimiento y ciclos económicos. Documento de Docencia D-5, ILADES--Georgetown University. [12.5]
\item Velde, F. R. (2013). Bitcoin: A Primer. \emph{Chicago Fed Letter}, 317, Federal Reserve Bank of Chicago. [8.1]
\item Walsh, C. E. (1995). Optimal Contracts for Central Bankers. \emph{American Economic Review}, 85(1). [9.5]
\item Walsh, C. E. (2002). Teaching Inflation Targeting: An Analysis for Intermediate Macro. \emph{Journal of Economic Education}, 33(4), 333--347. [13.1]
\item Werning, I. (2011). Managing a Liquidity Trap: Monetary and Fiscal Policy. NBER Working Paper 17344. [13.5]
\item Werning, I. (2015). Incomplete Markets and Aggregate Demand. NBER Working Paper 21448. [15.9]
\item Williamson, S. (2010). \emph{Macroeconomics}, 5.ª ed. Addison Wesley. [4.1]
\item Woodford, M. (2003). \emph{Interest and Prices: Foundations of a Theory of Monetary Policy}. Princeton University Press. [8.4, 8.5, 13.1]
\end{itemize}
}

\subsection*{Otras referencias con horas}
Lecturas que no proceden de los programas de la PUCP (líneas \emph{Artículos});
entre corchetes, el tema donde se usan.

\medskip
{\small
\begin{itemize}[leftmargin=1.4em,itemsep=1pt,label={}]
\item Auclert, A., B. Bardóczy, M. Rognlie y L. Straub (2021). Using the Sequence-Space Jacobian to Solve and Estimate Heterogeneous-Agent Models. \emph{Econometrica}, 89(5), 2375--2408. [15.8]
\item Céspedes, L. F., R. Chang y A. Velasco (2004). Balance Sheets and Exchange Rate Policy. \emph{American Economic Review}, 94(4), 1183--1193. [13.7]
\item Huggett, M. (1993). The Risk-Free Rate in Heterogeneous-Agent Incomplete-Insurance Economies. \emph{Journal of Economic Dynamics and Control}, 17(5--6), 953--969. [15.6]
\item Krusell, P. y A. A. Smith, Jr. (1998). Income and Wealth Heterogeneity in the Macroeconomy. \emph{Journal of Political Economy}, 106(5), 867--896. [15.7]
\item Levy Yeyati, E. (2006). Financial Dollarization: Evaluating the Consequences. \emph{Economic Policy}, 21(45), 61--118. [13.7]
\end{itemize}
}

%------------------------------------------------------------------------------
\newpage
\section*{Anexo G. Material opcional de consulta}
\addcontentsline{toc}{section}{Anexo G. Material opcional de consulta}

Todo lo de este anexo es opcional: no entra en el cómputo de horas ni modifica el
cronograma. Sirve para practicar, profundizar o ampliar el plan cuando el ritmo lo
permita. Los materiales de G.1 y G.2 aparecen además en las líneas
«Consulta:» de los temas indicados.

\subsection*{G.1 Práctica y ejercicios resueltos}
{\small\filasalternas
\noindent\begin{tabularx}{\textwidth}{>{\raggedright\arraybackslash}p{4.1cm}>{\raggedright\arraybackslash}X>{\raggedright\arraybackslash}p{3.2cm}}
\toprule
\rowcolor{tonoOscuro}\cab{Recurso} & \cab{Uso sugerido} & \cab{Temas} \\
\midrule
Romer, solucionario de \emph{Advanced Macroeconomics} & Problemas resueltos de una edición anterior (11 capítulos). Su numeración no coincide con la 5.ª edición: conviene buscar los problemas por tema y no por número. & Temas con lecturas de Romer \\[2pt]
Irigoyen, Rossi-Hansberg y Wright (2002), solucionario de Stokey, Lucas y Prescott & Ejercicios resueltos de \emph{Recursive Methods in Economic Dynamics}. Conviene intentar al menos los de los caps.~2 a 5 y 9 antes de mirar la solución. & 14.1, 14.2 \\[2pt]
Jiménez (2007), Tomo II: \emph{Ejercicios resueltos} & Ejercicios de nivel intermedio de los capítulos 3 a 8, 10 y 11. Quedan fuera los capítulos 1 y 2 (introducción y cuentas nacionales) y el 9 (balanza de pagos). & 1.1, 1.5, 2.1, 8.3, 8.4, 8.8, 9.2 y Bloque 12 (cap.~11) \\[2pt]
Prácticas UP (2024): dirigidas de Macroeconomía I, Universidad del Pacífico & Problemas con microfundamentos: modelos estáticos (secc.~1), dinámicos de economía cerrada (secc.~2), de economía abierta (secc.~3) y elección óptima del gobierno (secc.~4). & 14.3 (secc.~1); 2.1, 2.2 y 2.5 (secc.~2); 2.7 (secc.~4); la secc.~3 acompaña al Bloque~3 \\
\bottomrule
\end{tabularx}
}

\subsection*{G.2 Métodos computacionales}
{\small\filasalternas
\noindent\begin{tabularx}{\textwidth}{>{\raggedright\arraybackslash}p{4.1cm}>{\raggedright\arraybackslash}X>{\raggedright\arraybackslash}p{3.2cm}}
\toprule
\rowcolor{tonoOscuro}\cab{Recurso} & \cab{Uso sugerido} & \cab{Temas} \\
\midrule
Heer y Maußner (2005), \emph{Dynamic General Equilibrium Modelling} & Caps.~1 a 4 (algoritmos básicos, métodos lineal-cuadráticos y log-lineales, expectativas parametrizadas, métodos de proyección); caps.~5 a 7 (agentes heterogéneos y generaciones solapadas); caps.~8 y 9 (herramientas numéricas, cadenas de Markov y filtro HP). & 12.1 (secc.~9.4), 14.2, 14.5, 14.7, 15.6 \\[2pt]
Sánchez y Maraví (2009), notas de clase (UNMSM) & Solow, Ramsey, log-linealización y programación dinámica, con diagramas de flujo para programar. & 14.5, 14.7 \\
\bottomrule
\end{tabularx}
}

\subsection*{G.3 Módulo opcional de macroeconometría}
Sigue el temario del curso de extensión del BCRP de Eduardo Morón (2005), que
está en tu carpeta, con las notas de Sørensen (1996) como apoyo de series de
tiempo. Puede recorrerse en paralelo a los Bloques 8 y 12.

\medskip
{\small\filasalternas
\noindent\begin{tabularx}{\textwidth}{>{\raggedright\arraybackslash}p{4.1cm}>{\raggedright\arraybackslash}X>{\raggedright\arraybackslash}p{3.2cm}}
\toprule
\rowcolor{tonoOscuro}\cab{Tema} & \cab{Lecturas} & \cab{Relación con el plan} \\
\midrule
Series de tiempo: raíces unitarias, cointegración, ARCH & Sørensen (1996) & 8.2, 12.1 \\
VAR estructurales: identificación y estimación & Stock y Watson (2001); Blanchard y Quah (1989); Bernanke y Mihov (1998) & 8.2, 12.7 \\
Método generalizado de momentos & Hansen y West (2002); Davidson y MacKinnon (1999), cap.~9 & 12.3, 12.7 \\
Modelos de estado-espacio y filtro de Kalman & Elekdag (2001) & 12.11 \\
Modelos no lineales: umbral y transición suave & Escribano y Jordá (1997); van Dijk, Teräsvirta y Franses (2001) & Sin equivalente en el plan \\
Calibración de modelos & Material del tema 14.8 & 14.8 \\
\bottomrule
\end{tabularx}
}

\subsection*{G.4 Proyecto aplicado opcional: un modelo calibrado para el Perú}
Consolida los Bloques 12 a 14. Momento sugerido: después del Bloque 14 o en
paralelo al Bloque 15; no tiene horas asignadas.
\begin{enumerate}[leftmargin=1.6em,itemsep=2pt]
  \item \textbf{Datos.} Series trimestrales de producto, consumo, inversión y horas del
        BCRP (el complemento BCRPData de tu disco D facilita la descarga). Extraer el
        componente cíclico y sus momentos (tema 12.1; \GM{1.2}; Heer y Maußner,
        secc.~9.4).
  \item \textbf{Modelo.} El RBC básico del tema 12.5, calibrado con datos peruanos
        (\GM{3} y \GM{5}).
  \item \textbf{Solución y simulación} en Dynare (temas 14.5 y 14.8; \GM{2.2--2.4};
        \CSV{Apéndices B y C}).
  \item \textbf{Evaluación.} Comparar los momentos simulados con los observados (tema
        12.7) y, si se quiere, aplicar la contabilidad del ciclo (tema 12.8).
  \item \textbf{Extensión.} Un modelo neokeynesiano pequeño con regla de Taylor
        (temas 12.10, 13.1 y 13.2).
\end{enumerate}

\subsection*{G.5 Lecturas transversales y complementos}
\noindent\destacado{Historia del pensamiento.} Snowdon y Vane (2005), \emph{Modern
Macroeconomics}, se puede leer en paralelo como hilo histórico:
\begin{itemize}[leftmargin=1.4em,itemsep=1pt]
  \item Caps.~1 a 3 (Keynes frente a los clásicos; escuela keynesiana ortodoxa): Bloque 1.
  \item Cap.~4 (escuela monetarista): Bloque 8. Cap.~5 (nueva escuela clásica): Bloque 9.
  \item Cap.~6 (ciclo económico real): Bloque 12. Cap.~7 (nuevos keynesianos): Bloques 11 a 13.
  \item Cap.~10 (nueva macroeconomía política): temas 2.8, 9.4 y 9.5.
        Cap.~11 (renacimiento de la investigación sobre crecimiento): Bloques 4 a 6.
  \item Caps.~8, 9 y 12 (poskeynesianos, escuela austriaca y conclusiones): sin
        equivalente directo en el plan.
\end{itemize}

\medskip
\noindent\destacado{Macroeconomía internacional.} Complementan el Bloque 3, que no se
modifica:
\begin{itemize}[leftmargin=1.4em,itemsep=1pt]
  \item Obstfeld y Rogoff (1996): caps.~1 y 2 (comercio intertemporal y dinámica de
        la economía pequeña), tema 3.2; cap.~3 (ciclo de vida, política tributaria y
        cuenta corriente), tema 3.2 y Bloque 7; cap.~4 (tipo de cambio real y
        términos de intercambio), temas 3.3 y 3.4.
  \item Végh (2013): cap.~1 (modelo intertemporal básico), tema 3.2; caps.~3 y 4
        (precios intertemporales, bienes no transables y precios relativos), temas
        3.3 y 3.4.
  \item Uribe y Schmitt-Grohé (2015): caps.~1 a 5 (ciclos en economías abiertas y
        emergentes), tema 3.2 y tema 12.5; cap.~7 (transables, no transables y
        términos de intercambio), tema 3.4.
\end{itemize}
\noindent El sílabo de economía internacional desarrolla con horas todo este
material, junto con el comercio internacional.

\medskip
\noindent\destacado{Programas PUCP de crecimiento no integrados.} No se usaron sus
lecturas; sirven como guía para profundizar en los Bloques 4 a 6.
\begin{itemize}[leftmargin=1.4em,itemsep=1pt]
  \item 1ECO19, Teoría del Crecimiento (2026-2): unidades 1 y 2 (Solow, Ramsey, AK),
        temas 4.1 a 4.7 y 5.1; unidad 3 (Uzawa--Lucas, Romer, Aghion--Howitt), temas
        5.2, 5.4 y 5.5; unidad 4 (capital público, salud, intangibles), temas 4.8 y
        5.1; unidad 5 (Hsieh--Klenow, Hopenhayn), temas 5.5 y 6.2; unidad 6
        (Nelson--Phelps, Parente--Prescott, cambio estructural, automatización), temas
        5.2, 6.2 y 6.3; unidad 7 (innovación verde, clima, población), temas 4.9, 5.6
        y 6.5; unidad 8 (debates), Bloque 6.
  \item 1ECO85, Crecimiento y Desarrollo Económico (2025-1): el módulo 1 (Malthus,
        revolución industrial, modelo neoclásico, crecimiento endógeno) acompaña a
        los temas 4.2, 4.7, 5.1 y 6.5; los módulos 2 a 4 (economía política del
        desarrollo, desarrollo humano y pobreza) quedan fuera del plan.
\end{itemize}

\subsection*{G.6 Módulo opcional: macroeconomía e inteligencia artificial}
Reúne trabajos recientes que usan modelos de crecimiento y de macroeconomía para
evaluar el efecto de la inteligencia artificial. No tiene horas ni semanas
asignadas y no modifica el cronograma. Conviene abrirlo después del Bloque 5 (como
mínimo, hasta el tema 5.3); el último subgrupo, después del Bloque 14. Muchos
textos son documentos de trabajo, ensayos o notas docentes que aún no están en
revistas, de modo que su versión puede cambiar: conviene revisar la vigente antes
de citarlos. Se encontraron en publicaciones de economistas en X.

\tema{G.6.1 Automatización por tareas y eslabones débiles}
\lecturas{Steinsson y Hassanein (2026), notas docentes sobre automatización basada
en tareas \quad Jones y Tonetti (2026), «Past Automation and Future AI: How Weak
Links Tame the Growth Explosion»}
\nota{Punto de partida del módulo. Las notas de Steinsson y Hassanein están pensadas
para maestría o para un buen curso de crecimiento de grado: un continuo de tareas,
capital y trabajo como sustitutos perfectos dentro de cada tarea, y una elasticidad
$\sigma$ entre tareas. Con $\sigma<1$ el producto queda acotado por el trabajo
(«eslabones débiles»), y la automatización reduce la participación laboral. Jones y
Tonetti aplican la idea al crecimiento. Se apoya en los temas 4.1, 5.3, 5.6 y 6.3.}

\tema{G.6.2 ¿Cuánto crecimiento puede traer la IA?}
\lecturas{Moll y coautores (2026), «Will AI Soon Deliver Double-Digit Growth?»
\quad Korinek, Jones, Sacher, Cotter y McCrory (2026), «Economic Scenarios for
Transformative AI»}
\nota{El primero es un ensayo publicado en el blog de Alex Imas; falta verificar la
lista completa de autores. El segundo propone un marco para comparar escenarios a
2030: el PIB queda entre 1,6\,\% y 32\,\% sobre la trayectoria sin IA, y la
participación del trabajo baja de 60\,\% a entre 56\,\% y 45\,\%. Se apoya en los
temas 4.1, 5.7 y 10.1.}

\tema{G.6.3 Productividad, invención y ciencia}
\lecturas{Imas y Manyika (2026), «Bending the Curve of Discovery» \quad Agrawal,
Gans y Goldfarb, «Genius on Demand: The Value of Transformative Artificial
Intelligence», cap.~1 de Agrawal, Brynjolfsson y Korinek (eds.), \emph{The Economics
of Transformative AI}, con comentario de Garicano}
\nota{La IA como «invención de un método de invención» y el valor de una
inteligencia disponible a pedido. Se apoya en los temas 5.3 a 5.5. El volumen
completo (NBER y University of Chicago Press) sirve de guía del campo.}

\tema{G.6.4 Trabajo, salarios y distribución}
\lecturas{Acemoglu, Autor y Johnson (2026), «Building Pro-Worker Artificial
Intelligence», NBER w34854}
\nota{Dirección del cambio técnico y bienestar de los trabajadores. Continúa el
tema 5.6 y se relaciona con el Bloque 10.}

\tema{G.6.5 Tasas de interés, mercados y financiamiento de la inversión}
\lecturas{Chow, Halperin y Mazlish (2023), «Transformative AI, Existential Risk, and
Real Interest Rates» \quad Van Nieuwerburgh (2026), «Financing the AI Buildout»,
\emph{Brookings Papers on Economic Activity} y NBER w35865}
\nota{Con la ecuación de Euler del modelo de Ramsey (tema 4.7), más crecimiento
esperado implica tasas reales más altas. El primer texto usa esa relación para leer
qué esperan los mercados de la IA. El segundo estima una inversión de 10,3 billones
de dólares en 8 años y una facturación de 3,7 billones al año hacia 2032 para un
retorno de 10\,\% (cerca de 9\,\% del PIB). Se apoya en los temas 4.7, 4.8 y 8.8.}

\tema{G.6.6 La IA como herramienta de método}
\lecturas{Kase, Melosi y Rottner, «Estimating Nonlinear Heterogeneous Agent Models
with Neural Networks»}
\nota{Redes neuronales para estimar modelos con agentes heterogéneos. Complementa
los temas 14.7 y 15.9. No trata el efecto de la IA sobre la economía.}

\subsection*{Referencias del material opcional y de consulta}
{\small
\begin{itemize}[leftmargin=1.4em,itemsep=1pt,label={}]
\item Acemoglu, D., D. Autor y S. Johnson (2026). Building Pro-Worker Artificial Intelligence. NBER Working Paper 34854.
\item Agrawal, A., E. Brynjolfsson y A. Korinek (eds.) (2026). \emph{The Economics of Transformative AI}. NBER y University of Chicago Press. Incluye Agrawal, Gans y Goldfarb, «Genius on Demand: The Value of Transformative Artificial Intelligence» (cap.~1).
\item Armas, A., A. Ize y E. Levy Yeyati (eds.) (2006). \emph{Financial Dollarization: The Policy Agenda}. Palgrave Macmillan y Fondo Monetario Internacional.
\item Bernanke, B. y I. Mihov (1998). Measuring Monetary Policy. \emph{Quarterly Journal of Economics}, 113(3), 869--902.
\item Blanchard, O. J. y D. Quah (1989). The Dynamic Effects of Aggregate Demand and Supply Disturbances. \emph{American Economic Review}, 79(4), 655--673.
\item Chow, T., B. Halperin y J. Z. Mazlish (2023). Transformative AI, Existential Risk, and Real Interest Rates. Documento de trabajo.
\item Davidson, R. y J. MacKinnon (1999). The Generalized Method of Moments. Cap.~9 de \emph{Estimation and Inference in Econometrics}.
\item Elekdag, S. (2001). The Kalman Filter and Its Applications. Central Bank of Turkey.
\item Escribano, A. y O. Jordá (1997). Improved Testing and Specification of Smooth Transition Regression Models. University of California, Davis.
\item Finn, M. G. (2000). Perfect Competition and the Effects of Energy Price Increases on Economic Activity. \emph{Journal of Money, Credit and Banking}, 32(3), 400--416.
\item Hansen, B. y K. West (2002). Generalized Method of Moments and Macroeconomics. \emph{Journal of Business and Economic Statistics}.
\item Heer, B. y A. Maußner (2005). \emph{Dynamic General Equilibrium Modelling: Computational Methods and Applications}. Springer.
\item Imas, A. y J. Manyika (2026). Bending the Curve of Discovery. Ensayo para el Google DeepMind Institute.
\item Irigoyen, C., E. Rossi-Hansberg y M. L. J. Wright (2002). \emph{Solutions Manual for Recursive Methods in Economic Dynamics}. Harvard University Press.
\item Jiménez, F. (2007). \emph{Macroeconomía: enfoques y modelos. Tomo II: Ejercicios resueltos}, 2.ª ed. Fondo Editorial de la PUCP.
\item Jones, C. I. y C. Tonetti (2026). Past Automation and Future AI: How Weak Links Tame the Growth Explosion. Documento de trabajo, Stanford University.
\item Kaplan, G. y G. L. Violante (2018). Microeconomic Heterogeneity and Macroeconomic Shocks. \emph{Journal of Economic Perspectives}, 32(3), 167--194.
\item Kase, H., L. Melosi y M. Rottner. Estimating Nonlinear Heterogeneous Agent Models with Neural Networks. Documento de trabajo, Bank for International Settlements (WP 1241).
\item Kilian, L. (2008). The Economic Effects of Energy Price Shocks. \emph{Journal of Economic Literature}, 46(4), 871--909.
\item Kim, I.-M. y P. Loungani (1992). The Role of Energy in Real Business Cycle Models. \emph{Journal of Monetary Economics}, 29(2), 173--189.
\item Korinek, A., C. I. Jones, S. Sacher, T. Cotter y P. McCrory (2026). Economic Scenarios for Transformative AI. Anthropic.
\item Kydland, F. E. y E. C. Prescott (1977). Rules Rather than Discretion: The Inconsistency of Optimal Plans. \emph{Journal of Political Economy}, 85(3), 473--491.
\item Moll, B. y coautores (2026). Will AI Soon Deliver Double-Digit Growth? Ensayo en el blog de Alex Imas.
\item Morón, E. (2005). Macroeconometría. Programa del 52.º Curso de Extensión Universitaria, Banco Central de Reserva del Perú.
\item Obstfeld, M. y K. Rogoff (1996). \emph{Foundations of International Macroeconomics}. MIT Press.
\item Poole, W. (1968). Commercial Bank Reserve Management in a Stochastic Model: Implications for Monetary Policy. \emph{The Journal of Finance}, 23(5), 769--791.
\item Romer, D. Solucionario de \emph{Advanced Macroeconomics} (edición anterior, 11 capítulos).
\item Rotemberg, J. J. y M. Woodford (1996). Imperfect Competition and the Effects of Energy Price Increases on Economic Activity. \emph{Journal of Money, Credit and Banking}, 28(4), 550--577.
\item Sánchez, E. y C. Maraví (2009). Notas de clase de macroeconomía avanzada con métodos computacionales. IDES, Universidad Nacional Mayor de San Marcos.
\item Snowdon, B. y H. R. Vane (2005). \emph{Modern Macroeconomics: Its Origins, Development and Current State}. Edward Elgar.
\item Sørensen, B. E. (1996). Class Notes for Macroeconometrics (EC266). Brown University.
\item Steinsson, J. y A. Hassanein (2026). Notas docentes sobre automatización basada en tareas.
\item Stock, J. H. y M. W. Watson (2001). Vector Autoregressions. \emph{Journal of Economic Perspectives}, 15(4), 101--115.
\item Universidad del Pacífico (2024). Dirigidas para el curso de Macroeconomía I, semestre 2024-I (citadas como «Prácticas UP» en los temas).
\item Uribe, M. y S. Schmitt-Grohé (2015). \emph{Open Economy Macroeconomics}. Borrador del 17 de marzo de 2015, Columbia University.
\item van Dijk, D., T. Teräsvirta y P. H. Franses (2001). Smooth Transition Autoregressive Models: A Survey of Recent Developments. Econometric Institute Report 200, Erasmus University Rotterdam.
\item Van Nieuwerburgh, S. (2026). Financing the AI Buildout. \emph{Brookings Papers on Economic Activity}; NBER Working Paper 35865.
\item Végh, C. A. (2013). \emph{Open Economy Macroeconomics in Developing Countries}. MIT Press.
\end{itemize}
}

\vspace{1.5cm}
\franjas[2pt]
\begin{center}\small\color{tonoPizarra}
Plan de lectura de elaboración propia. Los títulos de los capítulos corresponden
a las obras citadas y se reproducen únicamente con fines de referencia
bibliográfica. La asignación de los capítulos del texto suplementario se basa en
sus títulos y en la descripción editorial, no en el manuscrito; la de Romer,
Campante et al., Acemoglu, Walsh, Galí y Galindo y Montecinos, en sus índices. Los temas y lecturas de la PUCP proceden de los
programas citados en el Anexo E.
\end{center}

\end{document}
